\documentclass[11pt]{article}

\usepackage[a4paper,margin=1in]{geometry}
\usepackage[T1]{fontenc}
\usepackage[utf8]{inputenc}
\usepackage{lmodern}
\usepackage{microtype}
\usepackage{amsmath,amssymb,amsthm,mathtools,bm}
\usepackage{booktabs}
\usepackage{enumitem}
\usepackage[numbers,sort&compress]{natbib}
\usepackage[colorlinks=true,linkcolor=blue,citecolor=blue,urlcolor=blue]{hyperref}
\usepackage[nameinlink,capitalize,noabbrev]{cleveref}

\allowdisplaybreaks

\newtheorem{assumption}{Assumption}
\newtheorem{definition}{Definition}
\newtheorem{theorem}{Theorem}
\newtheorem{proposition}{Proposition}
\newtheorem{corollary}{Corollary}
\newtheorem{lemma}{Lemma}
\newtheorem{remark}{Remark}

\crefname{assumption}{Assumption}{Assumptions}
\Crefname{assumption}{Assumption}{Assumptions}

\DeclareMathOperator{\Tr}{Tr}
\DeclareMathOperator{\diag}{diag}
\DeclareMathOperator{\KL}{KL}
\newcommand{\rhoid}{\rho_\star}
\newcommand{\rhop}{\rho_p}
\newcommand{\muid}{\mu_\star}
\newcommand{\MSE}{\operatorname{MSE}}
\newcommand{\Bias}{\operatorname{Bias}}
\newcommand{\Var}{\operatorname{Var}}
\newcommand{\Cov}{\operatorname{Cov}}
\newcommand{\E}{\mathbb{E}}
\newcommand{\Prob}{\mathbb{P}}
\newcommand{\R}{\mathbb{R}}
\newcommand{\N}{\mathbb{N}}
\newcommand{\1}{\mathbf{1}}
\newcommand{\OO}{\mathcal{O}}
\newcommand{\Deltaop}{\Delta}
\newcommand{\norm}[1]{\left\lVert #1 \right\rVert}
\newcommand{\abs}[1]{\left\lvert #1 \right\rvert}
\newcommand{\ip}[2]{\left\langle #1\middle|#2\right\rangle}
\newcommand{\bra}[1]{\left\langle #1\right|}
\newcommand{\ket}[1]{\left|#1\right\rangle}

\title{\bfseries Finite-Shot Operating Windows for\\ Virtual Distillation and Symmetry Verification}
\author{Anonymous Author}
\date{May 29, 2026}

\begin{document}
\maketitle


\begin{abstract}
Quantum error mitigation methods are often justified by infinite-shot bias reduction, while practical use is governed by finite sampling budgets and method-dependent resource costs.  This paper develops a theory-only operating-window analysis for virtual distillation (VD) and symmetry verification (SV).  The central result is a certified finite-shot mean-squared-error law for the VD quotient estimator \(\widehat\mu_{\mathrm{VD},M}=\widehat N_M/\widehat D_M\).  The law separates four effects: the population bias of \(\rho_p^M/\Tr(\rho_p^M)\), the dominant-eigenvector mismatch floor, the statistical \(B^{-1}\) bias of a random quotient, and denominator-driven variance inflation controlled by \(D_M=\Tr(\rho_p^M)\).  Unlike a purely asymptotic delta-method statement, the quotient law is accompanied by explicit Bernstein denominator concentration, a clipped-ratio regularization argument, and remainder constants expressed in bounded moments of the numerator-denominator one-shot estimator.  For VD implemented through independent controlled-derangement/Hadamard-test estimators, the numerator and denominator variances are computed explicitly, giving a concrete logarithmic copy-number scale.  For SV, the corresponding law is derived from a precise first-order projection formula: detectable noise is removed by postselection, while the undetectable component inside the symmetry sector determines the residual bias and the acceptance probability determines the sampling penalty.  These method-specific laws are converted into certified crossing rules: a leading indifference curve is accepted only when a remainder bound verifies that it lies in the regime where the local law is valid.  A closed one-qubit model gives finite-\(p\) two-sided boundary bounds, and a restricted Le Cam lower bound explains when finite sampling or non-identifiability imposes limits no selector can evade.  The paper also provides accepted-transcript efficiency statements, a certified selection trichotomy, explicit toy-model certificates, a generic three-level toy model with nonzero VD and SV first-order residual bias, an Efron--Stein total-variance certificate for the clipped VD quotient, an asymptotic second-order quotient-variance diagnostic, and a multiplicative implementation-noise baseline for controlled-derangement/Hadamard-test implementations.  No numerical or hardware experiments are claimed; the contribution is a mathematical finite-budget framework identifying when VD or SV is favored by certified operating windows and when the constants should be interpreted as conservative proof certificates rather than operational thresholds.
\end{abstract}

\section{Introduction}

Quantum error mitigation (QEM) tries to improve observable estimates on noisy devices without implementing full fault tolerance.  Standard techniques include zero-noise extrapolation, probabilistic error cancellation, symmetry verification, Clifford-data regression, and multi-copy purification methods such as virtual distillation and error suppression by derangements \citep{Temme2017,LiBenjamin2017,Endo2018,BonetMonroig2018,McArdle2019,Czarnik2021,Lowe2021,Huggins2021,Koczor2021PRX,Bultrini2023,Ijaz2024}.  A common issue is that lower bias is not free.  Mitigation can increase shot noise, circuit count, circuit width, calibration cost, or all of these.

This paper focuses on two methods with analyzable finite-shot structure.  Virtual distillation (VD) estimates an observable with respect to
\begin{equation}
    \frac{\rho_p^M}{\Tr(\rho_p^M)},
\end{equation}
where \(M\) noisy copies are combined \citep{Huggins2021,Koczor2021PRX}.  VD suppresses subdominant spectral components, but it cannot remove the mismatch between the dominant eigenvector of the noisy state and the ideal state \citep{Koczor2021NJP}.  Symmetry verification (SV) projects samples onto a valid symmetry sector and discards invalid samples \citep{BonetMonroig2018,McArdle2019,Kakkar2022}.  Its residual bias comes from symmetry-preserving errors that postselection cannot detect; its variance cost comes from the acceptance probability.  Cai's symmetry expansion framework already gives a broad bias-versus-sampling view for symmetry-based mitigation \citep{Cai2021}.  The present paper uses the narrower SV case as one component in an explicit operating-window theory.

The main mathematical issue is that VD is a quotient estimator.  The population target is
\begin{equation}
    \mu_{\mathrm{VD},M}(p)=\frac{N_M(p)}{D_M(p)}
    =\frac{\Tr(O\rho_p^M)}{\Tr(\rho_p^M)},
\end{equation}
while the finite-shot estimator is \(\widehat N_M/\widehat D_M\).  Therefore a VD MSE law must account not only for population bias and variance, but also for the statistical bias of a random quotient and for the probability that the estimated denominator approaches zero.  Prior VD analyses identify the population mechanism and noise floors \citep{Huggins2021,Koczor2021PRX,Koczor2021NJP}, while circuit-level work studies implementation noise and special cases \citep{Teo2022,Vikstal2024}.  The contribution here is to place these effects inside a finite-budget MSE boundary.

The paper makes seven claims, all conditional and theory-level.
\begin{enumerate}[leftmargin=2em]
    \item It derives the VD quotient MSE law with explicit \(B^{-1}\) quotient bias and denominator-driven variance.
    \item It connects the VD bias-floor coefficient \(\delta_{\mathrm{VD}}\) directly to the first-order noise generator \(\mathcal L\), the observable \(O\), and the ideal pure state \(\rho_\star\).
    \item It proves a denominator concentration certificate showing when the quotient expansion is trustworthy.
    \item It gives a precise first-order SV bias formula and separates variance-driven \(p/B\) dependence from quotient-bias terms, which vanish for Bernoulli postselection.
    \item It adds a validity certificate for indifference curves, so a predicted crossing is accepted only when it falls inside the local-law regime.
    \item It includes a restricted lower-bound sanity check: if two ideal targets induce statistically indistinguishable noisy transcripts at budget \(R\), no mitigation selector can estimate both below the corresponding two-point risk.
    \item It adds two structural organizing results that do not depend on numerical experiments: an accepted-transcript efficiency proposition for postselected laws and a selection trichotomy that classifies all two-method comparisons into uniform dominance or a certified crossing.
\end{enumerate}

No empirical validation is included.  The goal is to build the mathematical core before choosing QAOA instances, derangement circuits, noise models, or hardware backends.

\section{Problem setup and local laws}

Let \(\rho_\star\) be the ideal output state and let \(O=O^\dagger\) be a bounded observable.  The target is
\begin{equation}
    \mu_\star=\Tr(O\rho_\star).
\end{equation}
The noisy output state at physical noise scale \(p\geq0\) is \(\rho_p\).  We assume the local expansion
\begin{equation}
    \rho_p=\rho_\star+p\Delta+\OO(p^2),
    \qquad
    \Delta=\mathcal L(\rho_\star),
    \label{eq:noise_expansion}
\end{equation}
where \(\Delta=\Delta^\dagger\) and \(\Tr\Delta=0\).  The trace condition follows from \(\Tr\rho_p=1\).

A mitigation method \(m\) produces an estimator \(\widehat\mu_m\).  Its performance is
\begin{equation}
    \MSE_m(p,B)=\E[(\widehat\mu_m-\mu_\star)^2]
    =\Bias(\widehat\mu_m)^2+\Var(\widehat\mu_m),
\end{equation}
where \(B\) denotes the number of independent statistical samples for that method.

\begin{definition}[Certified local finite-shot law]
\label{def:certified_law}
A method \(m\) has a certified local finite-shot law on a parameter set \(\mathcal P\) if there exist functions \(b_m(p)\), \(v_m(p)\), \(c_m(p)\), constants \(C_m(p)\), and a threshold \(B_{\min,m}(p)\) such that, for every \(p\in\mathcal P\) and \(B\geq B_{\min,m}(p)\),
\begin{equation}
    \MSE_m(p,B)=b_m(p)^2+\frac{v_m(p)+2b_m(p)c_m(p)}{B}+\rho_m(p,B),
    \qquad
    \abs{\rho_m(p,B)}\leq \frac{C_m(p)}{B^2}.
    \label{eq:certified_law}
\end{equation}
We write \(\widetilde v_m(p)=v_m(p)+2b_m(p)c_m(p)\).  Some ratio estimators first produce a two-part bound
\begin{equation}
    \abs{\rho_m(p,B)}
    \leq \frac{C_m(p)}{B^2}+C_m^{\exp}(p)e^{-\gamma_m(p)B},
    \qquad \gamma_m(p)>0.
    \label{eq:certified_law_with_exp}
\end{equation}
Such a statement is converted into \cref{eq:certified_law} by increasing the threshold to any
\begin{equation}
    B\geq B_{\exp,m}(p)
    \quad\text{such that}\quad
    C_m^{\exp}(p)e^{-\gamma_m(p)B}\leq \frac{C_{m,\mathrm{abs}}(p)}{B^2},
    \label{eq:exponential_absorption_threshold}
\end{equation}
and replacing \(C_m\) by \(C_m+C_{m,\mathrm{abs}}\).  Thus exponential bad-denominator terms are not silently hidden in the local law; they are either displayed explicitly or absorbed after an explicit sample-size threshold.
\end{definition}

The cross-term \(2b_mc_m/B\) is retained because it is the leading interaction between population bias and statistical quotient bias.  When a method is a plain sample mean, \(c_m=0\).  When its variance coefficient depends smoothly on \(p\), the local expansion of \(\widetilde v_m(p)/B\) generally contains \(p/B\) terms.  Those terms should not be hidden inside \(\OO(p^2/B)\) unless the first derivative vanishes.

For unmitigated estimation,
\begin{equation}
    \mu_0(p)=\Tr(O\rho_p),
    \qquad
    b_0(p)=\mu_0(p)-\mu_\star=\beta_0p+\OO(p^2),
    \qquad
    \beta_0=\Tr(O\Delta).
    \label{eq:unmitigated_bias}
\end{equation}
If the one-shot variance is \(\sigma_0^2(p)\), then
\begin{equation}
    \MSE_0(p,B)=b_0(p)^2+\frac{\sigma_0^2(p)}{B}.
\end{equation}

\section{Virtual distillation: population law}

For \(M\geq2\), idealized VD targets
\begin{equation}
    \mu_{\mathrm{VD},M}(p)=\frac{N_M(p)}{D_M(p)},
    \qquad
    N_M(p)=\Tr(O\rho_p^M),
    \qquad
    D_M(p)=\Tr(\rho_p^M).
    \label{eq:vd_population_law}
\end{equation}
Assume the spectral decomposition
\begin{equation}
    \rho_p=\sum_{j\geq1}\lambda_j(p)\ket{\phi_j(p)}\bra{\phi_j(p)},
    \qquad
    \lambda_1(p)>\lambda_2(p)\geq\cdots\geq0.
\end{equation}
Then
\begin{equation}
    \mu_{\mathrm{VD},M}(p)=
    \frac{\sum_j\lambda_j(p)^M O_{jj}(p)}{\sum_j\lambda_j(p)^M},
    \qquad
    O_{jj}(p)=\bra{\phi_j(p)}O\ket{\phi_j(p)}.
\end{equation}
VD suppresses the subdominant eigenvalues, but it converges to the observable expectation in the dominant eigenvector:
\begin{equation}
    \lim_{M\to\infty}\mu_{\mathrm{VD},M}(p)=O_{11}(p).
\end{equation}
The infinite-copy floor is therefore
\begin{equation}
    b_{\mathrm{VD},\infty}(p)=O_{11}(p)-\mu_\star.
\end{equation}

\begin{lemma}[Spectral leakage]
\label{lem:spectral_leakage_v2}
Let
\begin{equation}
    r_M(p)=\sum_{j>1}\left(\frac{\lambda_j(p)}{\lambda_1(p)}\right)^M.
\end{equation}
Then
\begin{equation}
    \abs{\mu_{\mathrm{VD},M}(p)-O_{11}(p)}
    \leq \frac{2\norm{O}_\infty r_M(p)}{1+r_M(p)}
    \leq2\norm{O}_\infty r_M(p).
\end{equation}
\end{lemma}

\begin{proof}
Divide numerator and denominator by \(\lambda_1^M\).  The difference from \(O_{11}\) is a weighted average of \(O_{jj}-O_{11}\) over \(j>1\), divided by \(1+r_M\).  Since \(\abs{O_{jj}-O_{11}}\leq2\norm O_\infty\), the bound follows.
\end{proof}

\subsection{First-order coefficient in terms of the noise generator}

Assume now that \(\rho_\star=\ket{\psi_\star}\bra{\psi_\star}\) is pure.  Let
\begin{equation}
    Q_\star=I-\ket{\psi_\star}\bra{\psi_\star}.
\end{equation}
The eigenvalue \(1\) of \(\rho_\star\) is simple and separated from the zero eigenspace by a unit gap.  Standard non-degenerate eigenvector perturbation gives
\begin{equation}
    \ket{\phi_1(p)}=\ket{\psi_\star}+p\ket{\eta}+\OO(p^2),
    \qquad
    \ket{\eta}=Q_\star\Delta\ket{\psi_\star},
    \qquad
    \ip{\psi_\star}{\eta}=0.
    \label{eq:eta_generator}
\end{equation}

\begin{proposition}[VD floor coefficient]
\label{prop:vd_delta_generator}
Under \cref{eq:noise_expansion} and \cref{eq:eta_generator},
\begin{equation}
    b_{\mathrm{VD},\infty}(p)=\delta_{\mathrm{VD}}p+\OO(p^2),
\end{equation}
with
\begin{equation}
    \delta_{\mathrm{VD}}
    =2\operatorname{Re}\bra{\psi_\star}OQ_\star\Delta\ket{\psi_\star}
    =2\operatorname{Re}\bra{\psi_\star}O(I-\ket{\psi_\star}\bra{\psi_\star})\mathcal L(\rho_\star)\ket{\psi_\star}.
    \label{eq:delta_vd_explicit}
\end{equation}
\end{proposition}

\begin{proof}
By expansion of the dominant eigenvector,
\begin{align}
    O_{11}(p)
    &=\bra{\psi_\star+p\eta}O\ket{\psi_\star+p\eta}+\OO(p^2)\\
    &=\mu_\star+2p\operatorname{Re}\bra{\psi_\star}O\ket{\eta}+\OO(p^2).
\end{align}
Substitute \(\ket\eta=Q_\star\Delta\ket{\psi_\star}\).
\end{proof}

This formula is important because it prevents \(\delta_{\mathrm{VD}}\) from being a free symbol.  It is an observable-dependent projection of the first-order noise perturbation onto the orthogonal subspace of the ideal state.


\section{Virtual distillation: denominator concentration and quotient MSE}

Let \((X_{M,s},Y_{M,s})_{s=1}^B\) be independent copies of a one-sample pair satisfying
\begin{equation}
    \E X_M=N_M(p),
    \qquad
    \E Y_M=D_M(p),
\end{equation}
and define
\begin{equation}
    \widehat N_M=\frac1B\sum_{s=1}^B X_{M,s},
    \qquad
    \widehat D_M=\frac1B\sum_{s=1}^B Y_{M,s},
    \qquad
    \widehat\mu_{\mathrm{VD},M}=\frac{\widehat N_M}{\widehat D_M}.
\end{equation}
Write
\begin{equation}
    \Sigma_M(p)=
    \begin{pmatrix}
        \sigma_{N,M}^2(p) & \sigma_{ND,M}(p)\\
        \sigma_{ND,M}(p) & \sigma_{D,M}^2(p)
    \end{pmatrix}.
\end{equation}
The ratio is meaningful only when the denominator is statistically separated from zero.  This point is part of the operating-window theory, not a technical nuisance.

\begin{assumption}[Bounded one-shot VD estimator]
\label{ass:bounded_vd_pair_v4}
For fixed \((p,M)\), \(D_M(p)>0\), and the numerator-denominator one-shot pair satisfies
\begin{equation}
    \abs{X_M}\leq K_{N,M}(p),
    \qquad
    \abs{Y_M-D_M(p)}\leq K_{D,M}(p)
\end{equation}
almost surely.  Define centered variables \(\xi=X_M-N_M\) and \(\zeta=Y_M-D_M\), and bounded moment constants
\begin{equation}
    m_{ab,M}(p)=\E[\xi^a\zeta^b],
    \qquad
    \bar m_{ab,M}(p)=\E[\abs{\xi}^a\abs{\zeta}^b].
\end{equation}
\end{assumption}

\begin{lemma}[Bernstein denominator certificate]
\label{lem:denom_bernstein_v4}
Under \cref{ass:bounded_vd_pair_v4},
\begin{equation}
    \Prob\!\left(\widehat D_M<\frac{D_M}{2}\right)
    \leq
    \exp\!\left[-\frac{B D_M(p)^2}{8\sigma_{D,M}^2(p)+\frac{4}{3}K_{D,M}(p)D_M(p)}\right].
    \label{eq:bernstein_denominator_v4}
\end{equation}
Consequently, to make this probability at most \(\varepsilon\), it is sufficient that
\begin{equation}
    B\geq
    B_{\mathrm{den},M}(p,\varepsilon)
    :=
    \left(\frac{8\sigma_{D,M}^2(p)}{D_M(p)^2}
    +\frac{4K_{D,M}(p)}{3D_M(p)}\right)
    \log\!\frac1\varepsilon.
    \label{eq:denom_sample_threshold_v4}
\end{equation}
\end{lemma}

\begin{proof}
Apply Bernstein's inequality to \(B^{-1}\sum_s(Y_{M,s}-D_M)\) with threshold \(D_M/2\) \citep{Boucheron2013}.
\end{proof}

The ratio \(\sigma_{D,M}^2/D_M^2\) in \cref{eq:denom_sample_threshold_v4} is the denominator-driven inflation parameter.  If \(D_M\) is small, VD may still have a favorable population bias but require many samples before the ratio estimator is stable.

For a non-asymptotic statement define the clipped quotient
\begin{equation}
    \widehat\mu_{\mathrm{VD},M}^{\mathrm{clip}}
    =\frac{\widehat N_M}{\max\{\widehat D_M,D_M/2\}}.
    \label{eq:clipped_quotient_v4}
\end{equation}
It equals the ordinary quotient on the good event \(\widehat D_M\geq D_M/2\).  The clipping is a proof device; in practice it is also the natural regularization when a random denominator is used.


\begin{definition}[Explicit quotient-remainder constants]
\label{def:quotient_constants_v4}
For fixed \((p,M)\), put \(\mu_M=N_M/D_M\).  For sample means
\(\bar\xi=B^{-1}\sum_s \xi_s\) and \(\bar\zeta=B^{-1}\sum_s\zeta_s\), define finite moment envelopes
\begin{align}
    \mathfrak m_{12,M}(p)&=\sup_{B\ge1} B^2\abs{\E[\bar\xi\bar\zeta^2]},
    &
    \mathfrak m_{21,M}(p)&=\sup_{B\ge1} B^2\abs{\E[\bar\xi^2\bar\zeta]},\\
    \mathfrak m_{03,M}(p)&=\sup_{B\ge1} B^2\abs{\E[\bar\zeta^3]},
    &
    \mathfrak m_{13,M}(p)&=\sup_{B\ge1} B^2\E[\abs{\bar\xi}\abs{\bar\zeta}^3],\\
    \mathfrak m_{04,M}(p)&=\sup_{B\ge1} B^2\E[\abs{\bar\zeta}^4],
    &
    \mathfrak m_{40,M}(p)&=\sup_{B\ge1} B^2\E[\abs{\bar\xi}^4],\\
    \mathfrak m_{22,M}(p)&=\sup_{B\ge1} B^2\E[\bar\xi^2\bar\zeta^2].
\end{align}
These constants are finite under \cref{ass:bounded_vd_pair_v4}.  A fully explicit admissible choice follows from standard fourth-moment estimates for bounded centered averages; for example
\(\mathfrak m_{04,M}\le 3\sigma_{D,M}^4+K_{D,M}^4\),
\(\mathfrak m_{40,M}\le 3\sigma_{N,M}^4+16K_{N,M}^4\),
\(\mathfrak m_{22,M}\le (3\sigma_{N,M}^4+16K_{N,M}^4)^{1/2}(3\sigma_{D,M}^4+K_{D,M}^4)^{1/2}\), and analogous bounded third-moment estimates give finite \(\mathfrak m_{12,M},\mathfrak m_{21,M},\mathfrak m_{03,M}\).

For the variance certificate define the linear and quadratic quotient terms
\begin{equation}
    L_M=\frac{\bar\xi-\mu_M\bar\zeta}{D_M},
    \qquad
    Q_M=\frac{\mu_M\bar\zeta^2-\bar\xi\bar\zeta}{D_M^2},
\end{equation}
and let \(T_M\) denote the signed sum of all third-and-higher good-event terms in the Neumann expansion after \(Q_M\), including the signed fourth-order tail.  The remaining finite variance-tail envelope is
\begin{equation}
    \mathfrak t^{(V)}_M(p)
    :=\sup_{B\ge1}B^2
    \abs{2\Cov(L_M,T_M)+2\Cov(Q_M,T_M)+\Var(T_M)}.
    \label{eq:tail_variance_envelope_v14}
\end{equation}
It is finite under the boundedness assumption; if desired it can be expanded into sixth- and eighth-moment envelopes of \((\bar\xi,\bar\zeta)\).  Keeping it as \(\mathfrak t^{(V)}_M\) avoids hiding the \(\Cov(L_M,T_M)\) term inside an informal \(\OO(B^{-2})\).
Define
\begin{align}
    A_{E,M}(p)
    &=\frac{\mathfrak m_{12,M}(p)+\abs{\mu_M}\mathfrak m_{03,M}(p)}{D_M(p)^3}
      +\frac{\mathfrak m_{13,M}(p)}{D_M(p)^4}
      +\frac{2K_{N,M}(p)\mathfrak m_{04,M}(p)}{D_M(p)^5},
       \label{eq:AE_constant_v4}\\
    A_{Q,M}(p)
    &=\frac{2}{D_M(p)^4}\left[
      \abs{\mu_M}^2\mathfrak m_{04,M}(p)+\mathfrak m_{22,M}(p)
      \right],
      \label{eq:AQ_constant_v14}\\
    A_{LQ,M}(p)
    &=\frac{2}{D_M(p)^3}\left[
      2\abs{\mu_M}\mathfrak m_{12,M}(p)+\mathfrak m_{21,M}(p)
      +\abs{\mu_M}^2\mathfrak m_{03,M}(p)
      \right],
      \label{eq:ALQ_constant_v14}\\
    A_{V,M}(p)
    &=A_{Q,M}(p)+A_{LQ,M}(p)+\mathfrak t^{(V)}_M(p),
      \label{eq:AV_constant_v4}\\
    E_{\mathrm{bad},M}(p,B)
    &=\frac{16K_{N,M}(p)^2}{D_M(p)^2}
    \exp\!\left[-\frac{B D_M(p)^2}{8\sigma_{D,M}^2(p)+\frac{4}{3}K_{D,M}(p)D_M(p)}\right].
\end{align}
The constants are conservative, but they now account explicitly for the mixed covariance between the linear and quadratic quotient terms.  In particular, \cref{eq:ALQ_constant_v14} controls \(2\Cov(L_M,Q_M)\), the term that is not controlled by bounding \(\Var(Q_M)\) alone.
\end{definition}

\begin{remark}[Tightness of the certified constants]
\label{rem:constant_tightness_v7}
The constants in \cref{eq:AV_constant_v4} are proof certificates, not sharp estimates of the physical sampling penalty.  The largest inflation comes from bounding the square of the second-order remainder by the worst-case inequality \((a_1+\cdots+a_k)^2\le k\sum_i a_i^2\).  This Cauchy--Schwarz step is intentionally robust but can overestimate the true second moment by several orders of magnitude.  Consequently, the small numerical parameters used in the toy non-vacuity checks below should be read as artifacts of a loose variance certificate rather than physical thresholds.  The Efron--Stein certificate in \cref{lem:efron_stein_vd_v11} and the direct second-order calculation in \cref{prop:vd_second_order_variance_v11} are included precisely to separate what is needed for a rigorous certificate from what is expected to control the true sampling scale.
\end{remark}

\begin{theorem}[Non-asymptotic VD quotient law with explicit moment remainder]
\label{thm:vd_quotient_certified_v4}
Under \cref{ass:bounded_vd_pair_v4}, for every \(B\geq B_{\mathrm{den},M}(p,\varepsilon)\), the clipped VD estimator satisfies
\begin{align}
    \E[\widehat\mu_{\mathrm{VD},M}^{\mathrm{clip}}]
    &=\mu_{\mathrm{VD},M}(p)+\frac{c_{\mathrm{VD},M}(p)}{B}+r^{(E)}_{M}(p,B),
    \label{eq:vd_expectation_certified_v4}\\
    \Var(\widehat\mu_{\mathrm{VD},M}^{\mathrm{clip}})
    &=\frac{v_{\mathrm{VD},M}(p)}{B}+r^{(V)}_{M}(p,B),
    \label{eq:vd_variance_certified_v4}
\end{align}
where
\begin{align}
    c_{\mathrm{VD},M}(p)
    &=\frac{\mu_{\mathrm{VD},M}(p)\sigma_{D,M}^2(p)-\sigma_{ND,M}(p)}{D_M(p)^2},
    \label{eq:vd_c_v4}\\
    v_{\mathrm{VD},M}(p)
    &=\frac{\sigma_{N,M}^2(p)+\mu_{\mathrm{VD},M}(p)^2\sigma_{D,M}^2(p)
    -2\mu_{\mathrm{VD},M}(p)\sigma_{ND,M}(p)}{D_M(p)^2}
    \label{eq:vd_v_v4}\\
    &=\frac{\Var(X_M-\mu_{\mathrm{VD},M}(p)Y_M)}{D_M(p)^2}.
\end{align}
Moreover,
\begin{align}
    \abs{r^{(E)}_{M}(p,B)}
    &\leq \frac{A_{E,M}(p)}{B^2}+\frac{4K_{N,M}(p)}{D_M(p)}
    \exp\!\left[-\frac{B D_M(p)^2}{8\sigma_{D,M}^2(p)+\frac{4}{3}K_{D,M}(p)D_M(p)}\right],
    \label{eq:explicit_E_remainder_v4}\\
    \abs{r^{(V)}_{M}(p,B)}
    &\leq \frac{A_{V,M}(p)}{B^2}+E_{\mathrm{bad},M}(p,B).
    \label{eq:explicit_V_remainder_v4}
\end{align}
Therefore
\begin{equation}
    \MSE_{\mathrm{VD},M}(p,B)
    =b_{\mathrm{VD},M}(p)^2
    +\frac{v_{\mathrm{VD},M}(p)+2b_{\mathrm{VD},M}(p)c_{\mathrm{VD},M}(p)}{B}
    +\rho_{\mathrm{VD},M}(p,B),
    \label{eq:vd_mse_certified_v4}
\end{equation}
with
\begin{equation}
    \abs{\rho_{\mathrm{VD},M}(p,B)}
    \leq
    \frac{C_{\mathrm{VD},M}(p)}{B^2}
    +C'_{\mathrm{VD},M}(p)
    \exp\!\left[-\frac{B D_M(p)^2}{8\sigma_{D,M}^2(p)+\frac{4}{3}K_{D,M}(p)D_M(p)}\right],
    \label{eq:vd_mse_remainder_v4}
\end{equation}
where, with
\begin{equation}
    B_{E,M}=\frac{4K_{N,M}}{D_M},
    \qquad
    B_{V,M}=\frac{16K_{N,M}^2}{D_M^2},
\end{equation}
one admissible explicit choice is
\begin{align}
    C_{\mathrm{VD},M}
    &=A_{V,M}+c_{\mathrm{VD},M}^2
      +2\abs{b_{\mathrm{VD},M}}A_{E,M}
      +2\abs{c_{\mathrm{VD},M}}A_{E,M}
      +2A_{E,M}^2,\\
    C'_{\mathrm{VD},M}
    &=B_{V,M}+2\abs{b_{\mathrm{VD},M}}B_{E,M}
      +2\abs{c_{\mathrm{VD},M}}B_{E,M}+2B_{E,M}^2.
\end{align}
These constants are not optimized; they are written only to make the certificate reproducible.  In the strict sense of \cref{def:certified_law}, choose any \(B_{\exp,\mathrm{VD},M}\) satisfying \cref{eq:exponential_absorption_threshold} with
\begin{equation}
    \gamma_{\mathrm{VD},M}(p)=\frac{D_M(p)^2}{8\sigma_{D,M}^2(p)+\frac43 K_{D,M}(p)D_M(p)},
    \qquad
    C_{\mathrm{VD},M}^{\exp}(p)=C'_{\mathrm{VD},M}(p),
\end{equation}
and absorb the exponential contribution into the \(B^{-2}\) remainder whenever
\[
    B\geq\max\{B_{\mathrm{den},M},B_{\exp,\mathrm{VD},M}\}.
\]
\end{theorem}

\begin{proof}
Let \(G=\{\abs{\bar\zeta}\le D_M/2\}\).  On \(G\), clipping is inactive and \(\widehat D_M\ge D_M/2\).  The centered variables satisfy \(\abs{N_M+\bar\xi}=\abs{\widehat N_M}\le K_{N,M}\).  Since \(\abs{\bar\zeta/D_M}\le1/2\), the Neumann expansion is absolutely convergent and
\begin{equation}
    \frac{1}{D_M+\bar\zeta}
    =\frac1{D_M}\left(1-\frac{\bar\zeta}{D_M}
    +\frac{\bar\zeta^2}{D_M^2}-\frac{\bar\zeta^3}{D_M^3}\right)
    +\frac1{D_M}\frac{(\bar\zeta/D_M)^4}{1+\bar\zeta/D_M}.
\end{equation}
Multiplying by \(N_M+\bar\xi=\mu_MD_M+\bar\xi\) gives the term-by-term identity
\begin{align}
    \frac{N_M+\bar\xi}{D_M+\bar\zeta}
    &=\mu_M+\frac{\bar\xi-\mu_M\bar\zeta}{D_M}
      +\frac{\mu_M\bar\zeta^2-\bar\xi\bar\zeta}{D_M^2}
      +\frac{\bar\xi\bar\zeta^2-\mu_M\bar\zeta^3}{D_M^3}
      -\frac{\bar\xi\bar\zeta^3}{D_M^4}
      +\widetilde{\mathcal R}_4,
      \label{eq:explicit_neumann_terms_v5}
\end{align}
where
\begin{equation}
    \widetilde{\mathcal R}_4
    =\frac{N_M+\bar\xi}{D_M}\frac{(\bar\zeta/D_M)^4}{1+\bar\zeta/D_M}.
\end{equation}
The good-event tail is therefore explicitly bounded by
\begin{equation}
    \abs{\widetilde{\mathcal R}_4}
    \le \frac{2K_{N,M}\abs{\bar\zeta}^4}{D_M^5}.
    \label{eq:r4_tail_bound_v5}
\end{equation}
This is the missing geometric-tail control: the mixed cubic term \(-\bar\xi\bar\zeta^3/D_M^4\) is kept separate, and the remaining tail contains only powers \(\bar\zeta^4\) and higher.  Taking expectations in \cref{eq:explicit_neumann_terms_v5}, using \(\E\bar\xi=\E\bar\zeta=0\), \(\E\bar\zeta^2=\sigma_{D,M}^2/B\), \(\E\bar\xi\bar\zeta=\sigma_{ND,M}/B\), and the definitions of \(\mathfrak m_{12},\mathfrak m_{03},\mathfrak m_{13},\mathfrak m_{04}\), gives
\begin{equation}
    \left|\E\left[\widehat\mu^{\rm clip}_{\mathrm{VD},M}\mathbf 1_G\right]
    -\mu_M-\frac{c_{\mathrm{VD},M}}{B}\right|
    \le \frac{A_{E,M}}{B^2}.
\end{equation}
The displayed moment calculations use unrestricted moments.  Restricting them to \(G\) changes each bounded monomial expectation by at most its uniform bound times \(\Prob(G^c)\); these corrections are absorbed into the same exponential bad-event term.  On \(G^c\), the clipped estimator obeys \(\abs{\widehat\mu^{\rm clip}_{\mathrm{VD},M}}\le 2K_{N,M}/D_M\), and \cref{lem:denom_bernstein_v4} gives the exponential contribution in \cref{eq:explicit_E_remainder_v4}.  This proves the expectation claim.

For the variance, write the good-event expansion as
\begin{equation}
    \widehat\mu^{\rm clip}_{\mathrm{VD},M}
    =\mu_M+L_M+Q_M+T_M,
    \qquad
    L_M=\frac{\bar\xi-\mu_M\bar\zeta}{D_M},
    \qquad
    Q_M=\frac{\mu_M\bar\zeta^2-\bar\xi\bar\zeta}{D_M^2},
\end{equation}
where \(T_M\) collects the third-and-higher terms from \cref{eq:explicit_neumann_terms_v5,eq:r4_tail_bound_v5}.  Then \(B\Var(L_M)=v_{\mathrm{VD},M}\).  The quadratic variance term obeys
\begin{equation}
    B^2\Var(Q_M)
    \le B^2\E[Q_M^2]
    \le A_{Q,M}.
\end{equation}
The mixed linear-quadratic term is controlled explicitly by
\begin{align}
    2B^2\abs{\Cov(L_M,Q_M)}
    &\le \frac{2B^2}{D_M^3}
    \abs{\E\big[(\bar\xi-\mu_M\bar\zeta)(\mu_M\bar\zeta^2-\bar\xi\bar\zeta)\big]} \\
    &\le \frac{2}{D_M^3}
    \left[2\abs{\mu_M}\mathfrak m_{12,M}+\mathfrak m_{21,M}
    +\abs{\mu_M}^2\mathfrak m_{03,M}\right]
    =A_{LQ,M}.
\end{align}
Finally, the remaining third-and-higher contributions are exactly those collected in \(\mathfrak t^{(V)}_M\).  Therefore
\begin{equation}
    \abs{\Var(L_M+Q_M+T_M)-\Var(L_M)}
    \le \frac{A_{Q,M}+A_{LQ,M}+\mathfrak t^{(V)}_M}{B^2},
\end{equation}
which is \cref{eq:explicit_V_remainder_v4} on the good event; the bad event is again bounded by \(E_{\mathrm{bad},M}\).  This explicitly repairs the covariance term that would be missed by controlling only the square of the remainder.  The MSE expansion follows by combining the variance bound with
\begin{equation}
    \left(b_{\mathrm{VD},M}+\frac{c_{\mathrm{VD},M}}{B}+r^{(E)}_M\right)^2
    =b_{\mathrm{VD},M}^2+\frac{2b_{\mathrm{VD},M}c_{\mathrm{VD},M}}{B}+\OO(B^{-2}),
\end{equation}
where the \(\OO(B^{-2})\) term is replaced by the explicit bounds above in \cref{eq:vd_mse_remainder_v4}.
\end{proof}

\begin{remark}[Ordinary versus clipped quotient]
If the implementation guarantees inverse-moment control for \(\widehat D_M^{-1}\), the same leading expansion applies to the ordinary quotient.  Without such a guarantee, the clipped quotient is the mathematically safe estimator.  Since the two estimators coincide except on an exponentially small event once \cref{eq:denom_sample_threshold_v4} holds, their leading operating-window laws are identical.
\end{remark}


\subsection{Sharper variance diagnostics for the VD quotient}
\label{subsec:vd_sharper_variance_v11}

The certified law above is intentionally worst-case.  This subsection adds two more precise diagnostics.  The first is a fully non-asymptotic Efron--Stein certificate for the total variance of the clipped quotient.  It does not recover the exact coefficient \(v_{\mathrm{VD},M}\), but it avoids the large second-order Cauchy--Schwarz envelope used in \(A_{V,M}\).  The second is an asymptotic second-order variance calculation.  It explains which moment combination controls the true \(B^{-2}\) correction; it is not used as a replacement for the non-asymptotic certificate unless its displayed remainder is separately bounded.

\begin{lemma}[Efron--Stein total-variance certificate for clipped VD]
\label{lem:efron_stein_vd_v11}
Under \cref{ass:bounded_vd_pair_v4}, the clipped quotient in \cref{eq:clipped_quotient_v4} satisfies
\begin{equation}
    \Var\!\left(\widehat\mu_{\mathrm{VD},M}^{\mathrm{clip}}\right)
    \le
    \frac{8K_{N,M}^2\bigl(D_M+2K_{D,M}\bigr)^2}{D_M^4}\,\frac1B.
    \label{eq:efron_stein_vd_bound_v11}
\end{equation}
Consequently, for \(D_M\ge1/2\) and bounded Hadamard-test signs \(K_{N,M},K_{D,M}=\OO(1)\), the clipped VD quotient has a non-asymptotic total-variance certificate with a constant often of order \(10^2\)--\(10^3\) in the bounded-sign cases considered below, rather than the much larger \(B^{-2}\)-remainder envelope in \cref{eq:AV_constant_v4}.
\end{lemma}

\begin{proof}
Let \(f\) be the clipped quotient as a function of the \(B\) independent pairs \((X_s,Y_s)\).  Replace only the \(j\)-th pair and denote the resulting function by \(f^{(j)}\).  The sample numerator changes by at most \(2K_{N,M}/B\), while the sample denominator changes by at most \(2K_{D,M}/B\).  Since the clipped denominator is always at least \(D_M/2\), the map \((n,d)\mapsto n/\max\{d,D_M/2\}\) is Lipschitz with constants \(2/D_M\) in \(n\) and at most \(4K_{N,M}/D_M^2\) in \(d\), using \(\abs{\widehat N_M}\le K_{N,M}\).  Hence
\begin{equation}
    \abs{f-f^{(j)}}
    \le
    \frac{4K_{N,M}}{BD_M}
    +\frac{8K_{N,M}K_{D,M}}{BD_M^2}
    =
    \frac{4K_{N,M}(D_M+2K_{D,M})}{BD_M^2}.
\end{equation}
Efron--Stein gives
\begin{equation}
    \Var(f)\le \frac12\sum_{j=1}^B \E\big[(f-f^{(j)})^2\big]
    \le \frac{B}{2}\left(\frac{4K_{N,M}(D_M+2K_{D,M})}{BD_M^2}\right)^2,
\end{equation}
which is \cref{eq:efron_stein_vd_bound_v11}.
\end{proof}

\begin{proposition}[Second-order variance coefficient of the quotient]
\label{prop:vd_second_order_variance_v11}
Assume the centered one-shot variables have bounded sixth moments.  Let \(\mu_M=N_M/D_M\), \(\bar\xi=B^{-1}\sum_s\xi_s\), \(\bar\zeta=B^{-1}\sum_s\zeta_s\), and define the linear and quadratic quotient terms
\begin{align}
    L_M&=\frac{\bar\xi-\mu_M\bar\zeta}{D_M},\\
    Q_M&=\frac{\mu_M\bar\zeta^2-\bar\xi\bar\zeta}{D_M^2}.
\end{align}
Then, on the good-denominator event and up to an exponentially small clipping correction,
\begin{equation}
    \Var\!\left(\widehat\mu_{\mathrm{VD},M}^{\mathrm{clip}}\right)
    =\frac{v_{\mathrm{VD},M}}{B}
     +\frac{W_{\mathrm{VD},M}^{(0)}}{B^2}
     +\OO(B^{-3})+\OO(e^{-\gamma_M B}),
    \label{eq:second_order_variance_expansion_v11}
\end{equation}
where the clean second-order coefficient is
\begin{align}
    W_{\mathrm{VD},M}^{(0)}
    &=\Gamma_{\mathrm{VD},M}
      +\frac{1}{D_M^4}\Big[
      2\mu_M^2\sigma_{D,M}^4
      +\sigma_{N,M}^2\sigma_{D,M}^2
      +\sigma_{ND,M}^2
      -4\mu_M\sigma_{D,M}^2\sigma_{ND,M}
      \Big],                                      \label{eq:second_order_W0_v13}\\
    2B^2\Cov(L_M,Q_M)
    &=\Gamma_{\mathrm{VD},M}+\OO(B^{-1}),          \label{eq:third_moment_gamma_v13}\\
    B^2\Var(Q_M)
    &=\frac{1}{D_M^4}\Big[
      2\mu_M^2\sigma_{D,M}^4
      +\sigma_{N,M}^2\sigma_{D,M}^2
      +\sigma_{ND,M}^2
      -4\mu_M\sigma_{D,M}^2\sigma_{ND,M}
      \Big]+\OO(B^{-1}).
      \label{eq:second_order_Q_variance_v11}
\end{align}
The term \(\Gamma_{\mathrm{VD},M}\) is the third-moment contribution to the quotient variance.  In symmetric Hadamard-test implementations it often vanishes or is lower order, but no such cancellation is assumed in the non-asymptotic certificate.
\end{proposition}

\begin{proof}
Use the Neumann expansion of \((D_M+\bar\zeta)^{-1}\) on the event \(\abs{\bar\zeta}\le D_M/2\):
\begin{equation}
    \frac{N_M+\bar\xi}{D_M+\bar\zeta}
    =\mu_M+L_M+Q_M+\OO\!\left(\abs{\bar\xi}\abs{\bar\zeta}^2+\abs{\bar\zeta}^3\right).
\end{equation}
The leading term satisfies \(B\Var(L_M)=v_{\mathrm{VD},M}\).  For centered i.i.d.\ sample averages,
\begin{align}
    \Var(\bar\zeta^2)&=\frac{2\sigma_{D,M}^4}{B^2}+\OO(B^{-3}),\\
    \Var(\bar\xi\bar\zeta)&=\frac{\sigma_{N,M}^2\sigma_{D,M}^2+\sigma_{ND,M}^2}{B^2}+\OO(B^{-3}),\\
    \Cov(\bar\zeta^2,\bar\xi\bar\zeta)&=\frac{2\sigma_{D,M}^2\sigma_{ND,M}}{B^2}+\OO(B^{-3}).
\end{align}
Substituting these three identities into \(Q_M=(\mu_M\bar\zeta^2-\bar\xi\bar\zeta)/D_M^2\) gives \cref{eq:second_order_Q_variance_v11}.  The remaining covariance \(\Cov(L_M,Q_M)\) contains third moments and is kept explicit.  Cubic and higher Neumann terms are \(\OO(B^{-3})\) in variance under the sixth-moment assumption, while the complement of the good event contributes only the exponential clipping term already controlled in \cref{lem:denom_bernstein_v4}.
\end{proof}

\begin{remark}[How this changes the interpretation of the toy certificates]
\label{rem:sharp_variance_interpretation_v11}
\Cref{lem:efron_stein_vd_v11,prop:vd_second_order_variance_v11} do not change the leading operating-window law.  They explain why the explicit constants used in the certified toy examples are pessimistic: \(A_{V,M}\) is a worst-case certificate for a second-order remainder, whereas the actual \(B^{-2}\) variance correction is controlled by the moment combination in \cref{eq:second_order_Q_variance_v11}.  Thus the earlier extreme numerical choices demonstrate formal non-vacuity of the certificate; they should not be read as the physical scale at which VD or SV becomes useful.
\end{remark}






\section{Copy-number choice for VD}

The copy number \(M\) is not a free improvement knob.  Increasing \(M\) suppresses subdominant spectral components, but it can also reduce \(D_M=\Tr(\rho_p^M)\), which increases the denominator-driven variance.

\begin{lemma}[Denominator spectral bounds]
\label{lem:denom_spectral_bounds_v4}
Let \(r_M(p)=\sum_{j>1}(\lambda_j/\lambda_1)^M\).  Then
\begin{equation}
    D_M(p)=\lambda_1(p)^M\bigl(1+r_M(p)\bigr),
    \qquad
    \lambda_1(p)^M\leq D_M(p)\leq \lambda_1(p)^M(1+r_M(p)).
    \label{eq:dm_spectral_bounds_v4}
\end{equation}
\end{lemma}

\begin{proof}
Use \(D_M=\sum_j\lambda_j^M\) and factor out \(\lambda_1^M\).
\end{proof}

\subsection{Concrete independent derangement/Hadamard-test implementation}

To make the variance constants non-aspirational, fix the following idealized implementation.  Let \(U_M\) be the cyclic derangement acting on \(M\) copies.  The denominator \(D_M=\Tr(\rho_p^M)\) is estimated by a controlled-\(U_M\) Hadamard test.  For a Pauli observable \(O\), the numerator \(N_M=\Tr(O\rho_p^M)\) is estimated by an independent controlled observable-derangement Hadamard test.  Each test returns a binary outcome in \(\{\pm1\}\).  We use \(B\) denominator shots and \(B\) numerator shots; hence the physical resource is proportional to \(R_M=2\kappa_M B\), where \(\kappa_M\) is the circuit resource per Hadamard test.

For this implementation,
\begin{equation}
    X_M\in\{\pm1\},
    \qquad
    Y_M\in\{\pm1\},
    \qquad
    \E X_M=N_M,
    \qquad
    \E Y_M=D_M,
\end{equation}
and the two streams are independent.  Therefore
\begin{equation}
    \sigma_{N,M}^2=1-N_M^2,
    \qquad
    \sigma_{D,M}^2=1-D_M^2,
    \qquad
    \sigma_{ND,M}=0.
    \label{eq:hadamard_variances_v4}
\end{equation}
The VD quotient constants become
\begin{align}
    c_{\mathrm{VD},M}^{\mathrm{Had}}(p)
    &=\frac{\mu_{\mathrm{VD},M}(p)\bigl(1-D_M(p)^2\bigr)}{D_M(p)^2},
    \label{eq:hadamard_c_v4}\\
    v_{\mathrm{VD},M}^{\mathrm{Had}}(p)
    &=\frac{1-N_M(p)^2+\mu_{\mathrm{VD},M}(p)^2\bigl(1-D_M(p)^2\bigr)}{D_M(p)^2}.
    \label{eq:hadamard_v_v4}
\end{align}
For Pauli \(O\) with \(\|O\|_\infty\le1\), positivity of \(\rho_p^M\) gives \(\abs{N_M}\le \Tr(\rho_p^M)=D_M\), hence \(\abs{\mu_{\mathrm{VD},M}}\le1\).  The denominator certificate is explicit:
\begin{equation}
    B\geq
    \left(
    \frac{8(1-D_M^2)}{D_M^2}+\frac{8}{3D_M}
    \right)\log\!\frac1\varepsilon
    \label{eq:hadamard_Bmin_v4}
\end{equation}
is sufficient for \(\Prob(\widehat D_M<D_M/2)\leq\varepsilon\).

\begin{proposition}[Implementation-level VD variance inflation]
\label{prop:implementation_variance_inflation_v4}
For the independent Hadamard-test implementation above,
\begin{equation}
    v_{\mathrm{VD},M}^{\mathrm{Had}}(p)
    \leq \frac{2}{D_M(p)^2}
    \leq 2\lambda_1(p)^{-2M}.
    \label{eq:implementation_upper_v4}
\end{equation}
If, in addition, \(1-N_M(p)^2\geq v_->0\) on the considered parameter range, then
\begin{equation}
    v_{\mathrm{VD},M}^{\mathrm{Had}}(p)\geq v_-\lambda_1(p)^{-2M}(1+r_M(p))^{-2}.
    \label{eq:implementation_lower_v4}
\end{equation}
\end{proposition}

\begin{proof}
The upper bound follows from \cref{eq:hadamard_v_v4} and \(D_M\geq\lambda_1^M\).  The lower bound uses the numerator-variance term \((1-N_M^2)/D_M^2\) and \(D_M=\lambda_1^M(1+r_M)\).
\end{proof}

\begin{proposition}[Multiplicative implementation-noise shift of the VD law]
\label{prop:impl_noise_shift_v11}
\textnormal{\textit{Model-dependent refinement.}}
\textnormal{\textit{This proposition is not a compiled-circuit noise theorem.}}
Suppose the controlled-derangement/Hadamard tests used to estimate the population quantities are followed by effective sign contractions
\begin{equation}
    \widetilde N_M=\eta_{N,M}N_M,
    \qquad
    \widetilde D_M=\eta_{D,M}D_M,
    \qquad
    \eta_{N,M},\eta_{D,M}>0,
\end{equation}
with \(\theta_M^{\rm impl}:=\eta_{N,M}/\eta_{D,M}\).  Then the implemented population quotient is
\begin{equation}
    \mu_{\mathrm{VD},M}^{\rm impl}=\theta_M^{\rm impl}\frac{N_M}{D_M},
\end{equation}
so the implemented population bias is
\begin{equation}
    b_{\mathrm{VD},M}^{\rm impl}
    =\left(\theta_M^{\rm impl}-1\right)\frac{N_M}{D_M}+b_{\mathrm{VD},M}.
\end{equation}
If the same contraction affects numerator and denominator, \(\theta_M^{\rm impl}=1\), the population quotient is unchanged and implementation noise appears first as sampling inflation.  In the independent sign model with one-shot variables in \(\{\pm1\}\),
\begin{equation}
    v_{\mathrm{VD},M}^{\rm impl}
    =\frac{1-\eta_{N,M}^2N_M^2+\big(\theta_M^{\rm impl}N_M/D_M\big)^2(1-\eta_{D,M}^2D_M^2)}{\eta_{D,M}^2D_M^2}.
    \label{eq:impl_noise_vd_variance_v11}
\end{equation}
Consequently, in the common-contraction case \(\eta_{N,M}=\eta_{D,M}=\eta_M\), the population quotient and hence the bias gap are unchanged.  The denominator part of the sampling coefficient carries an \(\eta_M^{-2}\)-type inflation, but the full coefficient is not generally a pure multiplicative rescaling of the ideal one because the Bernoulli variances are also replaced by \(1-\eta_M^2N_M^2\) and \(1-\eta_M^2D_M^2\).
\end{proposition}

\begin{proof}
The first two displays follow by taking the ratio of the contracted means.  For the variance coefficient, substitute the contracted means into the quotient variance formula \cref{eq:vd_v_v4} with independent numerator and denominator signs, so the covariance term is zero.  The denominator in \cref{eq:vd_v_v4} becomes \((\eta_{D,M}D_M)^2\), giving \cref{eq:impl_noise_vd_variance_v11}.  If \(\eta_{N,M}=\eta_{D,M}=\eta_M\), then the population quotient and hence the bias gap are unchanged.  The denominator in the quotient-variance formula contributes an \(\eta_M^{-2}\)-type inflation, but \cref{eq:impl_noise_vd_variance_v11} also replaces the Bernoulli variance factors by \(1-\eta_M^2N_M^2\) and \(1-\eta_M^2D_M^2\).  Thus the leading crossing budget should be computed from the displayed coefficient rather than by applying a universal scalar multiplier to the ideal variance.
\end{proof}

\begin{remark}[Interpretation]
\label{rem:impl_noise_interpretation_v11}
This model records the algebraic effect of multiplicative readout or controlled-test contractions on a VD quotient.  It is deliberately narrower than a circuit-level noise analysis: a compiled derangement implementation may create correlated numerator-denominator errors, coherent control errors, or nonmultiplicative distortions.  The proposition should therefore be read as a falsifiable baseline for the experimental section: common multiplicative test noise cancels in the population ratio but modifies the sampling coefficient through \cref{eq:impl_noise_vd_variance_v11}; asymmetric test noise creates an additional first-order bias.
\end{remark}


\begin{theorem}[Logarithmic copy-number scale for implemented VD]
\label{thm:mstar_log_scale_v4}
Fix \(p\) and suppose that, for \(M\) in an admissible range before the dominant-eigenvector floor takes over,
\begin{equation}
    A_-^2q^{2M}
    \leq
    \bigl(b_{\mathrm{VD},M}(p)-b_{\mathrm{VD},\infty}(p)\bigr)^2
    \leq
    A_+^2q^{2M},
    \qquad
    q=\frac{\lambda_2(p)}{\lambda_1(p)}<1, \qquad 0<\lambda_2(p)<\lambda_1(p),
    \label{eq:bias_two_sided_mstar_v4}
\end{equation}
and the Hadamard-test variance has the two-sided scaling in \cref{prop:implementation_variance_inflation_v4}.  Then every minimizer of the leading implemented VD objective
\begin{equation}
    \mathcal E_M(B)=A^2q^{2M}+\frac{V\lambda_1^{-2M}}{B}
\end{equation}
lies within an \(\OO(1)\) additive distance of
\begin{equation}
    M_B=\frac{\log B}{2\abs{\log\lambda_2(p)}}.
    \label{eq:mstar_log_lambda2_v4}
\end{equation}
Since \(\lambda_1(p)=1+\OO(p)\) in the local pure-state regime,
\begin{equation}
    M_B=\Theta\!\left(\frac{\log B}{\abs{\log\!\bigl(\lambda_2(p)/\lambda_1(p)\bigr)}}\right)
\end{equation}
up to local \(\lambda_1\)-dependent constants.  Once \(\abs{b_{\mathrm{VD},\infty}}\) is comparable to the leakage term, the optimal copy number saturates because additional copies increase variance without reducing the floor.
\end{theorem}

\begin{proof}
Balance the decreasing leakage term against the increasing denominator-driven variance term:
\begin{equation}
    q^{2M}\asymp B^{-1}\lambda_1^{-2M}.
\end{equation}
Equivalently, \((q\lambda_1)^{2M}\asymp B^{-1}\).  Since \(q\lambda_1=\lambda_2\), \cref{eq:mstar_log_lambda2_v4} follows.  The two-sided constants in \cref{eq:bias_two_sided_mstar_v4} and \cref{prop:implementation_variance_inflation_v4} shift the minimizer only by an additive constant independent of \(B\).  Integer rounding also changes \(M\) by at most a constant.
\end{proof}

\begin{remark}
The exact logarithmic denominator contains \(\abs{\log\lambda_2}\), because the variance inflation is driven by \(D_M\sim\lambda_1^M\) while leakage is driven by \((\lambda_2/\lambda_1)^M\).  In the small-noise regime \(\lambda_1\approx1\), this is equivalent to the more informal \(\log B/\abs{\log\!\bigl(\lambda_2/\lambda_1\bigr)}\) scaling.  This distinction should be kept in the paper because it identifies which part of VD controls the cost.
\end{remark}

\section{Symmetry verification}

Let \(P\) be the projector onto the valid symmetry sector and assume
\begin{equation}
    P\rho_\star P=\rho_\star.
\end{equation}
For a symmetry-compatible observable, SV estimates
\begin{equation}
    \mu_{\mathrm{SV}}(p)=\frac{\Tr(OP\rho_pP)}{a(p)},
    \qquad
    a(p)=\Tr(P\rho_p).
    \label{eq:sv_population_v2}
\end{equation}
The first-order projected state is
\begin{equation}
    \frac{P\rho_pP}{\Tr(P\rho_p)}
    =\rho_\star+p\left(P\Delta P-\rho_\star\Tr(P\Delta P)\right)+\OO(p^2).
    \label{eq:sv_projected_state_expansion}
\end{equation}

A precise detectable-undetectable split can be defined by
\begin{equation}
    \Delta_{\mathrm{undet}}=P\Delta P,
    \qquad
    P\Delta_{\mathrm{det}}P=0,
    \qquad
    \Delta=\Delta_{\mathrm{undet}}+\Delta_{\mathrm{det}}
\end{equation}
with \(\Delta_{\mathrm{det}}=\Delta-P\Delta P\) for this first-order analysis.  This is a definition of the split, not a convention.  The detectable part is exactly the part that has no component inside the postselected block.

\begin{proposition}[SV first-order bias]
\label{prop:sv_delta_precise}
Under \cref{eq:noise_expansion} and \(P\rho_\star P=\rho_\star\),
\begin{equation}
    a(p)=1+p\Tr(P\Delta)+\OO(p^2),
\end{equation}
and
\begin{equation}
    b_{\mathrm{SV}}(p)=\mu_{\mathrm{SV}}(p)-\mu_\star
    =\delta_{\mathrm{SV}}p+\OO(p^2),
\end{equation}
where
\begin{equation}
    \delta_{\mathrm{SV}}
    =\Tr(OP\Delta P)-\mu_\star\Tr(P\Delta P).
    \label{eq:delta_sv_precise}
\end{equation}
If one writes \(a(p)=1-\alpha_{\mathrm{det}}p+\OO(p^2)\), then \(\alpha_{\mathrm{det}}=-\Tr(P\Delta)\).
\end{proposition}

\begin{proof}
Use \(P\rho_pP=\rho_\star+pP\Delta P+\OO(p^2)\) and
\begin{equation}
    \frac{1}{1+p\Tr(P\Delta P)+\OO(p^2)}
    =1-p\Tr(P\Delta P)+\OO(p^2).
\end{equation}
Substitute into \(\Tr(O\,P\rho_pP)/\Tr(P\rho_p)\) and subtract \(\mu_\star\).
\end{proof}

Let \(A_s\in\{0,1\}\) be the acceptance indicator and \(W_s\) the measured value when accepted.  Put \(Z_s=A_sW_s\).  Then
\begin{equation}
    \widehat\mu_{\mathrm{SV}}=\frac{\overline Z}{\overline A},
    \qquad
    \E A=a(p),
    \qquad
    \E Z=a(p)\mu_{\mathrm{SV}}(p).
\end{equation}


\begin{theorem}[Certified SV postselection quotient law]
\label{thm:sv_quotient_law_v2}
Assume \(a(p)>0\), \(A\in\{0,1\}\), and \(\abs{W}\le K_W\) on accepted shots.  Put \(Z=AW\), \(\mu_{\mathrm{SV}}=\E Z/a\), and define centered variables
\begin{equation}
    \xi_{\mathrm{SV}}=Z-a\mu_{\mathrm{SV}},
    \qquad
    \zeta_{\mathrm{SV}}=A-a.
\end{equation}
Let \(\widehat\mu_{\mathrm{SV}}^{\rm clip}=\overline Z/\max\{\overline A,a/2\}\).  Then
\begin{equation}
    \Prob\left(\overline A<a/2\right)
    \le \exp\left[-\frac{B a(p)}{8}\right].
    \label{eq:sv_acceptance_bernstein_v5}
\end{equation}
Moreover,
\begin{align}
    \E[\widehat\mu_{\mathrm{SV}}^{\rm clip}]
    &=\mu_{\mathrm{SV}}(p)+\frac{c_{\mathrm{SV}}(p)}{B}+r^{(E)}_{\mathrm{SV}}(p,B),\\
    \Var(\widehat\mu_{\mathrm{SV}}^{\rm clip})
    &=\frac{v_{\mathrm{SV}}(p)}{B}+r^{(V)}_{\mathrm{SV}}(p,B),
\end{align}
where
\begin{align}
    c_{\mathrm{SV}}(p)
    &=\frac{\mu_{\mathrm{SV}}(p)\Var(A)-\Cov(Z,A)}{a(p)^2}=0,
    \label{eq:sv_quotient_bias_zero_v14}\\
    v_{\mathrm{SV}}(p)
    &=\frac{\Var(Z-\mu_{\mathrm{SV}}(p)A)}{a(p)^2}.
\end{align}
The equality \(c_{\mathrm{SV}}(p)=0\) is exact, because \(Z=AW\) and hence
\(\Cov(Z,A)=\E[ZA]-\E Z\,\E A=a\mu_{\mathrm{SV}}(1-a)=\mu_{\mathrm{SV}}\Var(A)\).
There are explicit bounded-moment constants \(A_{E,\mathrm{SV}}(p)\), \(A_{V,\mathrm{SV}}(p)\) obtained from \cref{def:quotient_constants_v4} by the replacement
\begin{equation}
    (N_M,D_M,X_M,Y_M,K_{N,M},K_{D,M})
    \mapsto
    (a\mu_{\mathrm{SV}},a,Z,A,K_W,1).
\end{equation}
Thus
\begin{align}
    \abs{r^{(E)}_{\mathrm{SV}}(p,B)}
    &\le \frac{A_{E,\mathrm{SV}}(p)}{B^2}+\frac{4K_W}{a(p)}e^{-Ba(p)/8},\\
    \abs{r^{(V)}_{\mathrm{SV}}(p,B)}
    &\le \frac{A_{V,\mathrm{SV}}(p)}{B^2}+\frac{16K_W^2}{a(p)^2}e^{-Ba(p)/8}.
\end{align}
Consequently
\begin{equation}
    \MSE_{\mathrm{SV}}(p,B)
    =b_{\mathrm{SV}}(p)^2+
    \frac{v_{\mathrm{SV}}(p)}{B}
    +\rho_{\mathrm{SV}}(p,B),
    \label{eq:sv_certified_mse_v5}
\end{equation}
with \(\abs{\rho_{\mathrm{SV}}(p,B)}\le C_{\mathrm{SV}}(p)B^{-2}+C'_{\mathrm{SV}}(p)e^{-Ba(p)/8}\).  If \(\sigma_{\mathrm{acc}}^2(p)=\Var(W\mid A=1)\), then
\begin{equation}
    v_{\mathrm{SV}}(p)=\frac{\sigma_{\mathrm{acc}}^2(p)}{a(p)}.
\end{equation}
Therefore, if \(\sigma_{\mathrm{acc}}^2(p)=\sigma_\star^2+\sigma_1p+\OO(p^2)\) and \(a(p)=1-\alpha_{\mathrm{det}}p+\OO(p^2)\), then
\begin{align}
    \MSE_{\mathrm{SV}}(p,B)
    &=\delta_{\mathrm{SV}}^2p^2
    +\frac{\sigma_\star^2+(\sigma_1+\alpha_{\mathrm{det}}\sigma_\star^2)p}{B}
    +\OO(p^2/B)+\mathcal R_{\mathrm{SV}}(p,B),
\end{align}
where no separate \(\OO(p/B)\) quotient-bias cross term appears: for Bernoulli postselection the coefficient \(c_{\mathrm{SV}}\) vanishes identically by \cref{eq:sv_quotient_bias_zero_v14}.  The certified theorem-level statement is \(\abs{\mathcal R_{\mathrm{SV}}(p,B)}\le C_{\mathrm{SV}}(p)B^{-2}+C'_{\mathrm{SV}}(p)e^{-Ba(p)/8}\).
\end{theorem}

\begin{proof}
The acceptance denominator is Bernoulli.  The multiplicative Chernoff bound gives \(\Prob(\overline A\le (1-\delta)a)\le e^{-Ba\delta^2/2}\); taking \(\delta=1/2\) gives \cref{eq:sv_acceptance_bernstein_v5}.  Conditional on the good event \(\abs{\overline A-a}\le a/2\), apply the same explicit Neumann expansion used in \cref{thm:vd_quotient_certified_v4} to \(g(z,a)=z/a\).  Because \(A\in\{0,1\}\) and \(\abs Z\le K_W\), the moment envelopes are finite and the bad-event contribution is exponentially small.  The variance identity follows from
\begin{equation}
    Z-\mu_{\mathrm{SV}}A=A(W-\mu_{\mathrm{SV}}),
\end{equation}
so \(\Var(Z-\mu_{\mathrm{SV}}A)=a(p)\sigma_{\mathrm{acc}}^2(p)\).  The displayed local expansion follows from Taylor expansion of \(1/a(p)\) and from \cref{prop:sv_delta_precise}.
\end{proof}

\begin{remark}[Position relative to symmetry expansion]
This SV section should be presented as derivative in mechanism from symmetry-based mitigation theory, especially \citet{Cai2021}.  The added value here is the operating-window boundary and the compatibility with the VD quotient law, not a new symmetry estimator.
\end{remark}

\section{PEC and CDR in the selector}

The rigorous selector should include only methods whose local constants are specified.  PEC can be inserted cleanly under standard first-order cancellation assumptions.  Suppose the PEC estimator is unbiased through first order, so
\begin{equation}
    b_{\mathrm{PEC}}(p)=\OO(p^2),
\end{equation}
and let \(\gamma(p)=1+\gamma_1p+\OO(p^2)\) denote the quasiprobability one-norm overhead.  Then a local PEC law has the form
\begin{equation}
    \MSE_{\mathrm{PEC}}(p,B)=\OO(p^4)+\frac{\gamma(p)^2\sigma_{\mathrm{PEC}}^2(p)}{B}+\OO(B^{-2}).
\end{equation}
Thus PEC fits the selector with \(s_{\mathrm{PEC}}\approx\kappa_{\mathrm{PEC}}\gamma^2\sigma_{\mathrm{PEC}}^2\), where \(\kappa_{\mathrm{PEC}}\) accounts for resource normalization.

CDR is different.  Its bias and variance depend on the training distribution, regression model, target observable, and allocation of shots between training and target circuits \citep{Czarnik2021,Lowe2021}.  Therefore CDR should not be included in the rigorous lower envelope unless a separate local law is proven for the chosen training model.  Here, CDR is treated as an admissible separate component, not a theorem-level member of the selector.

\section{Certified operating-window crossings}

Let \(R\) be a common physical resource budget.  If method \(m\) consumes \(\kappa_m(p)\) resource units per effective statistical sample, then \(B_m=R/\kappa_m(p)\).  A certified law becomes
\begin{equation}
    \MSE_m(p,R)=b_m(p)^2+\frac{s_m(p)}{R}+\rho_m(p,R),
    \qquad
    \abs{\rho_m(p,R)}\leq\frac{C_m(p)}{R^2}
\end{equation}
for \(R\geq R_{\min,m}(p)\), where \(s_m=\kappa_m\widetilde v_m\).


\begin{remark}[Resource normalization]
\label{rem:resource_normalization_v13}
Here \(R\) denotes the physical resource being charged in the comparison.  If an estimator uses an effective statistical sample count \(B_m\), all streams, copies, and circuit calls required to produce those samples must be included in \(\kappa_m\).  In particular, the independent VD Hadamard-test implementation of \cref{prop:implementation_variance_inflation_v4} uses one numerator stream and one denominator stream, so a paired VD sample budget \(B\) corresponds to physical cost \(R_{\mathrm{VD},M}=2\kappa_M B\).  Equivalently, when \(R\) counts physical Hadamard-test calls, \(\kappa_{\mathrm{VD},M}=2\kappa_M\).  Toy-model formulas below state explicitly when \(R\) is a paired-shot budget and when it is a physical call budget.
\end{remark}


\begin{theorem}[Certified crossing perturbation]
\label{thm:certified_crossing}
Consider two methods \(i,j\) at fixed \(p\).  Suppose
\begin{equation}
    \MSE_i=b_i^2+\frac{s_i}{R}+\rho_i(R),
    \qquad
    \MSE_j=b_j^2+\frac{s_j}{R}+\rho_j(R),
    \qquad
    \abs{\rho_m(R)}\leq\frac{C_m}{R^2}
\end{equation}
for \(R\geq R_{\min}\), and suppose the remainders are continuous.  Let
\begin{equation}
    g=b_j^2-b_i^2>0,
    \qquad
    \Delta s=s_i-s_j>0,
    \qquad
    R_0=\frac{\Delta s}{g},
    \qquad
    C=C_i+C_j.
\end{equation}
If \(R_0(1-\eta)\geq R_{\min}\) and
\begin{equation}
    \eta=\frac{4C}{gR_0^2}<\frac12,
    \label{eq:eta_crossing_condition}
\end{equation}
then there exists at least one true crossing \(R^\star\) in
\begin{equation}
    R^\star\in[R_0(1-\eta),R_0(1+\eta)].
    \label{eq:crossing_interval}
\end{equation}
In particular every certified crossing in this interval obeys
\begin{equation}
    \abs{R^\star-R_0}\leq \frac{4C}{gR_0}.
    \label{eq:crossing_shift_bound}
\end{equation}
If, in addition, \(\rho_i-\rho_j\) is differentiable and
\begin{equation}
    \sup_{R\in[R_0(1-\eta),R_0(1+\eta)]}
    \abs{(\rho_i-\rho_j)'(R)}
    <\frac{\Delta s}{R_0^2(1+\eta)^2},
    \label{eq:crossing_monotonicity_condition_v5}
\end{equation}
then the crossing in the interval is unique.
\end{theorem}

\begin{proof}
The leading difference is
\begin{equation}
    F_0(R)=\left(b_i^2-b_j^2\right)+\frac{s_i-s_j}{R}
    =-g+\frac{\Delta s}{R},
\end{equation}
with root \(R_0\).  At \(R_+=R_0(1+\eta)\),
\begin{equation}
    F_0(R_+)=-g\frac{\eta}{1+\eta},
\end{equation}
and at \(R_-=R_0(1-\eta)\),
\begin{equation}
    F_0(R_-)=g\frac{\eta}{1-\eta}.
\end{equation}
The remainder difference has magnitude at most \(C/R^2\).  Since \(\eta=4C/(gR_0^2)\) and \(\eta<1/2\), the signs at \(R_+\) and \(R_-\) are preserved.  Continuity then gives existence by the intermediate value theorem.  If \cref{eq:crossing_monotonicity_condition_v5} holds, then
\begin{equation}
    F'(R)=-\frac{\Delta s}{R^2}+(\rho_i-\rho_j)'(R)<0
\end{equation}
throughout the interval, so the crossing there is unique.
\end{proof}

The leading indifference formula
\begin{equation}
    R_{i\leftrightarrow j}(p)=\frac{s_i(p)-s_j(p)}{b_j(p)^2-b_i(p)^2}
    \label{eq:leading_crossing_formula}
\end{equation}
is therefore not accepted automatically.  It must pass two checks: \(R_{i\leftrightarrow j}\geq R_{\min}\) and \(4C/[gR_{i\leftrightarrow j}^2]\ll1\).  This is the self-consistency condition missing from a purely formal crossing curve.


\begin{theorem}[Certified selection trichotomy at fixed \(p\)]
\label{thm:selection_trichotomy_v10}
At fixed \(p\), suppose two methods \(i,j\) admit certified local laws
\begin{equation}
    \MSE_m(R)=b_m^2+\frac{s_m}{R}+\rho_m(R),
    \qquad
    \abs{\rho_m(R)}\le \frac{C_m}{R^2},
    \qquad R\ge R_{\min},
\end{equation}
with continuous remainders.  Let
\begin{equation}
    g=b_j^2-b_i^2,
    \qquad
    \Delta s=s_i-s_j,
    \qquad
    C=C_i+C_j.
\end{equation}
Then exactly one of the following leading-order regimes applies, with the indicated certified interpretation.
\begin{enumerate}[leftmargin=2.2em]
\item[\textnormal{(0)}] \textbf{Degenerate leading tie.}  If \(g=0\) and \(\Delta s=0\), the two leading local laws are identical.  No leading-order dominance or crossing is certified; the comparison depends on the signed remainders, higher-order coefficients, or a sharper local expansion.
\item[\textnormal{(I)}] \textbf{Uniform dominance of \(i\).}  If \(g>0\) and \(\Delta s\le0\), then method \(i\) has strictly smaller leading bias and no larger sampling coefficient, and \(i\) wins for all
\begin{equation}
    R\ge \max\left\{R_{\min},\sqrt{\frac{C}{g}}\right\}.
\end{equation}
If \(g=0\) and \(\Delta s<0\), then \(i\) has equal leading bias and strictly smaller sampling coefficient, and \(i\) wins for all
\begin{equation}
    R\ge \max\left\{R_{\min},\frac{C}{\abs{\Delta s}}\right\}.
\end{equation}
\item[\textnormal{(II)}] \textbf{Uniform dominance of \(j\).}  If \(g<0\) and \(\Delta s\ge0\), or if \(g=0\) and \(\Delta s>0\), the symmetric statement holds after swapping \(i\) and \(j\).
\item[\textnormal{(III)}] \textbf{Genuine tradeoff with a certified crossing.}  If \(g\Delta s>0\), set
\begin{equation}
    R_0=\frac{\Delta s}{g}>0,
    \qquad
    \eta=\frac{4C}{\abs g R_0^2}.
\end{equation}
If \(\eta<1/2\) and \(R_0(1-\eta)\ge R_{\min}\), then at least one true crossing lies in
\begin{equation}
    [R_0(1-\eta),R_0(1+\eta)].
\end{equation}
If the monotonicity condition \cref{eq:crossing_monotonicity_condition_v5} also holds, that crossing is unique; the lower-bias method wins above it and the lower-sampling method wins below it.
\end{enumerate}
\end{theorem}

\begin{proof}
Write
\begin{equation}
    \Phi(R)=\MSE_i(R)-\MSE_j(R)=-g+\frac{\Delta s}{R}+\rho_i(R)-\rho_j(R),
    \qquad
    \abs{\rho_i(R)-\rho_j(R)}\le\frac{C}{R^2}.
\end{equation}
Case (0) is immediate from \(g=\Delta s=0\).  Cases (I) and (II) follow because the leading comparison has the same sign for all large enough \(R\), and the \(C/R^2\) remainder cannot overturn it once \(R\) exceeds the displayed threshold.  In case (III), the leading root is \(R_0=\Delta s/g\).  After relabelling if needed so that \(g>0\) and \(\Delta s>0\), the hypotheses reduce exactly to \cref{thm:certified_crossing}.  That theorem gives existence in the certified interval and uniqueness under \cref{eq:crossing_monotonicity_condition_v5}.
\end{proof}

\begin{corollary}[Local location of operating windows]
\label{cor:window_location_v10}
Suppose
\begin{equation}
    b_m(p)=\delta_m p+\OO(p^2),
    \qquad
    s_m(p)=\sigma_m^2+s_m^{(1)}p+\OO(p^2).
\end{equation}
If \(\sigma_i^2\ne\sigma_j^2\), a genuine bias--variance tradeoff has a leading window at scale \(R_0=\Theta(p^{-2})\).  If \(\sigma_i^2=\sigma_j^2\) but \(s_i^{(1)}\ne s_j^{(1)}\), the leading window is instead at scale \(R_0=\Theta(p^{-1})\).  Thus the \(1/p\) crossings in the closed toy models below are not accidental; they arise because the two compared laws have matched leading shot cost and differ first at order \(p\) in their sampling coefficients.
\end{corollary}

\begin{proof}
Substitute the displayed expansions into \(g=b_j^2-b_i^2\) and \(\Delta s=s_i-s_j\).  If the constant term of \(\Delta s\) is nonzero, then \(R_0=\Delta s/g=\Theta(p^{-2})\).  If that constant term cancels, then \(\Delta s=\Theta(p)\) while \(g=\Theta(p^2)\), giving \(R_0=\Theta(p^{-1})\).
\end{proof}

\begin{corollary}[Uniform finite-interval sandwich]
\label{cor:finite_interval_sandwich}
Suppose on \(p\in[0,p_{\max}]\) one has computable bounds
\begin{equation}
    0<G_-(p)\leq b_j(p)^2-b_i(p)^2\leq G_+(p),
    \qquad
    0<S_-(p)\leq s_i(p)-s_j(p)\leq S_+(p).
\end{equation}
Then the leading crossing satisfies
\begin{equation}
    \frac{S_-(p)}{G_+(p)}\leq R_{i\leftrightarrow j}(p)\leq\frac{S_+(p)}{G_-(p)}.
\end{equation}
If the certified-remainder condition of \cref{thm:certified_crossing} holds uniformly on that interval, this becomes a two-sided finite-\(p\) operating-window certificate, not merely a Taylor expansion at \(p=0\).

\end{corollary}

\section{Restricted lower-bound sanity check}
\label{sec:restricted_lower_bound}

The selector constructed above compares concrete mitigation procedures.  It does not by itself prove that the selected method is globally optimal among all conceivable estimators.  A full optimality theorem would require specifying what information about the noise channel, ideal state family, measurements, and adaptive strategies is available.  This section adds a deliberately restricted lower bound.  Its role is modest but important: it identifies when an operating window is limited by statistical distinguishability, not merely by the chosen menu of methods.

Consider a local two-point model indexed by \(\theta\in\{+,-\}\).  The ideal values are
\begin{equation}
    \mu_\theta^\star=\Tr(O\rho_\theta^\star),
    \qquad
    \abs{\mu_+^\star-\mu_-^\star}=2h.
\end{equation}
At physical noise level \(p\), a resource budget \(R\) and a fixed admissible measurement/mitigation protocol induce transcript distributions \(\mathbb P_+^{(R)}\) and \(\mathbb P_-^{(R)}\).  These transcript distributions include all random circuits, accepted or rejected shots, auxiliary calibration samples, and classical postprocessing randomness used by the protocol.

\begin{theorem}[Two-point finite-budget lower bound]
\label{thm:two_point_lower_bound}
For any estimator \(\widehat\mu\) of the ideal value based on the transcript,
\begin{equation}
    \sup_{\theta\in\{+,-\}}
    \mathbb E_\theta\left[(\widehat\mu-\mu_\theta^\star)^2\right]
    \geq
    \frac{h^2}{2}\left(1-\operatorname{TV}\left(\mathbb P_+^{(R)},\mathbb P_-^{(R)}\right)\right),
    \label{eq:tv_lower_bound}
\end{equation}
where \(\operatorname{TV}\) denotes total variation distance.  Consequently, if
\begin{equation}
    \operatorname{KL}\left(\mathbb P_+^{(R)}\,\|\,\mathbb P_-^{(R)}\right)\leq R I_p h^2,
    \label{eq:local_kl_condition}
\end{equation}
then
\begin{equation}
    \sup_{\theta}
    \mathbb E_\theta\left[(\widehat\mu-\mu_\theta^\star)^2\right]
    \geq
    \frac{h^2}{2}\left(1-\sqrt{\frac{R I_p h^2}{2}}\right)_+.
    \label{eq:kl_lower_bound}
\end{equation}
Choosing \(h=(2RI_p)^{-1/2}\), whenever this perturbation lies inside the local model, gives the sampling lower bound
\begin{equation}
    \sup_\theta
    \mathbb E_\theta\left[(\widehat\mu-\mu_\theta^\star)^2\right]
    \geq
    \frac{1}{8RI_p}.
    \label{eq:one_over_r_lower_bound}
\end{equation}
\end{theorem}

\begin{proof}
If an estimator has squared error less than \(h^2\) under both hypotheses, thresholding it halfway between \(\mu_-^\star\) and \(\mu_+^\star\) distinguishes the two hypotheses.  Le Cam's two-point testing argument gives \cref{eq:tv_lower_bound}.  Pinsker's inequality gives
\begin{equation}
    \operatorname{TV}(\mathbb P_+^{(R)},\mathbb P_-^{(R)})
    \leq
    \sqrt{\frac12\operatorname{KL}(\mathbb P_+^{(R)}\|\mathbb P_-^{(R)})},
\end{equation}
which yields \cref{eq:kl_lower_bound}.  The choice of \(h\) gives \cref{eq:one_over_r_lower_bound}.
\end{proof}


\begin{proposition}[Explicit \(I_p\) for a one-qubit Bernoulli transcript]
\label{prop:explicit_ip_bernoulli_v5}
Consider the restricted local model
\begin{equation}
    \rho_t^\star=\frac{I+tX}{2},
    \qquad
    \rho_{p,t}=\frac{I+(1-p)tX}{2},
    \qquad
    \mu_t^\star=t,
\end{equation}
with \(t\in[-t_0,t_0]\) and one resource unit equal to one noisy \(X\)-measurement.  For the two hypotheses \(t=\pm h\), the one-shot transcript distributions are Bernoulli distributions with means \(\pm(1-p)h\) on \(\{\pm1\}\), and for \(\abs{h}\le[2(1-p)]^{-1}\),
\begin{equation}
    \KL(P_{+h}\|P_{-h})
    \le 4(1-p)^2h^2.
\end{equation}
Thus \cref{eq:local_kl_condition} holds with the explicit conservative per-resource information constant
\begin{equation}
    I_p^{\rm Bern}=4(1-p)^2.
    \label{eq:explicit_ip_bernoulli_v5}
\end{equation}
Consequently every estimator based on \(R\) such shots obeys the concrete lower bound
\begin{equation}
    \sup_{t\in\{+h,-h\}}
    \E_t[(\widehat\mu-t)^2]
    \ge \frac{1}{32R(1-p)^2},
\end{equation}
provided the optimizing \(h=[2RI_p^{\rm Bern}]^{-1/2}\) lies in the local interval.
\end{proposition}

\begin{proof}
For \(Y\in\{\pm1\}\), write
\begin{equation}
    q_+=\frac{1+a}{2},\qquad q_-=\frac{1-a}{2},\qquad a=(1-p)h.
\end{equation}
In this symmetric two-point model \(q_-=1-q_+\), which is the only reason the following simplification is valid.  The Bernoulli KL is
\begin{align}
    \KL(P_{+h}\|P_{-h})
    &=q_+\log\frac{q_+}{q_-}+(1-q_+)\log\frac{1-q_+}{1-q_-}\\
    &=(2q_+-1)\log\frac{q_+}{1-q_+}
      =a\log\frac{1+a}{1-a}.
      \label{eq:bernoulli_kl_corrected_v9}
\end{align}
For \(\abs a\le1/2\), \(\left|\log[(1+a)/(1-a)]\right|\le4\abs a\), hence \(a\log[(1+a)/(1-a)]\le4a^2\) and the KL is at most \(4a^2\).  This gives the conservative finite-interval value \(I_p^{\rm Bern}=4(1-p)^2\).  The same expression also gives the sharper local constant: since \(\log[(1+a)/(1-a)]=2a+O(a^3)\),
\begin{equation}
    \KL(P_{+h}\|P_{-h})=2(1-p)^2h^2+O(h^4),
\end{equation}
so the local two-point/Fisher normalization is
\begin{equation}
    I_{p,\mathrm{loc}}^{\rm Bern}=2(1-p)^2.
    \label{eq:explicit_ip_fisher_local_v8}
\end{equation}
Equivalently, using \(\left|\log[(1+a)/(1-a)]\right|\le2\abs a/(1-a^2)\) gives \(I_p^{\rm Bern}=2(1-p)^2/(1-a^2)\), which tends to \cref{eq:explicit_ip_fisher_local_v8} as \(a\to0\).  Below we keep the conservative finite-interval value when a uniform certificate is needed and separately record the sharper local benchmark when comparing constants.
\end{proof}

\begin{corollary}[Structural non-identifiability floor]
\label{cor:nonidentifiability_floor}
If two distinct ideal values satisfy \(\abs{\mu_+^\star-\mu_-^\star}=2h\) but induce the same noisy transcript distribution for every budget \(R\), then every estimator obeys
\begin{equation}
    \sup_{\theta\in\{+,-\}}
    \mathbb E_\theta\left[(\widehat\mu-\mu_\theta^\star)^2\right]
    \geq \frac{h^2}{2}.
\end{equation}
Thus a budget-independent floor is unavoidable whenever the ideal functional is not identifiable from the noisy experiment and side information allowed to the protocol.
\end{corollary}

\begin{remark}[How this lower bound should be used]
\label{rem:lower_bound_role}
This is not a universal no-go theorem for QEM\@.  It is a restricted benchmark for a declared statistical model.  Its value is diagnostic.  If a VD/SV/PEC operating window lies far above \cref{eq:one_over_r_lower_bound}, the method menu or constants may be improvable.  If it approaches the two-point lower bound inside the same model, the selector is close to statistically optimal in that restricted sense.  The corollary also clarifies why bias floors matter: some floors are not artifacts of a mitigation method but consequences of information that the noisy transcript does not contain.
\end{remark}



\begin{proposition}[Accepted-transcript efficiency of the selected law]
\label{prop:accepted_transcript_efficiency_v10}
Fix \(p\) and suppose method \(m\) operates in a bias-negligible regime, meaning
\(b_m(p)=o(R^{-1/2})\) on the budget range under consideration, so that
\begin{equation}
    \MSE_m(p,R)=\frac{s_m(p)}{R}+\OO(R^{-2}).
\end{equation}
Assume the per-resource transcript actually used by method \(m\) is a regular local statistical experiment with Fisher information \(J_m(p)\) for the corresponding population value.  Then
\begin{equation}
    s_m(p)\ge \frac{1}{J_m(p)}.
    \label{eq:transcript_crb_v10}
\end{equation}
Equality holds when the method's estimator is locally efficient for that transcript.  In the Pauli/Bernoulli accepted transcript used in the SV toy models, where each accepted observation has conditional variance \(\sigma_{\rm acc}^2(p)\) for the local mean parameter and acceptance probability \(a(p)\),
\begin{equation}
    J_{\rm SV}(p)=\frac{a(p)}{\sigma_{\rm acc}^2(p)},
    \qquad
    s_{\rm SV}(p)=\frac{\sigma_{\rm acc}^2(p)}{a(p)}=\frac{1}{J_{\rm SV}(p)}.
    \label{eq:sv_fisher_pauli_corrected_v13}
\end{equation}
For a non-Bernoulli or non-efficient accepted transcript, \cref{eq:sv_fisher_pauli_corrected_v13} should be read as an additional local-efficiency assumption, not as a universal identity.  Thus the constant-factor gaps in the two-point Le Cam comparisons below are artifacts of conservative two-point testing only after the transcript model has been fixed.
\end{proposition}

\begin{proof}
Over \(R\) resource units the transcript has Fisher information \(R J_m(p)\).  The local Cramer--Rao inequality gives \(\Var(\widehat\mu_m)\ge 1/[R J_m(p)]\) for locally unbiased regular estimators, yielding \cref{eq:transcript_crb_v10}.  For the Pauli/Bernoulli SV transcript, the number of accepted samples is \(N_{\rm acc}=\sum_{r=1}^R A_r\) with \(A_r\sim\mathrm{Bernoulli}(a(p))\).  The same multiplicative Chernoff bound used in \cref{thm:sv_quotient_law_v2} gives, for every \(0<\delta<1\),
\begin{equation}
    \Pr\!\left(\abs{N_{\rm acc}-a(p)R}\ge \delta a(p)R\right)
    \le 2\exp[-a(p)R\delta^2/3].
\end{equation}
Thus \(N_{\rm acc}=a(p)R+O_{\mathbb P}(\sqrt{R})\).  Conditional on acceptance, the accepted Pauli/mean transcript has local information \(1/\sigma_{\rm acc}^2(p)\) per accepted observation, and the sample mean over accepted observations is locally efficient in this one-dimensional mean experiment.  Hence the accepted transcript has Fisher information \(a(p)R/\sigma_{\rm acc}^2(p)+o(R)\), and its leading variance coefficient is \(\sigma_{\rm acc}^2(p)/a(p)\).  This proves \cref{eq:sv_fisher_pauli_corrected_v13} in the stated Pauli/Bernoulli setting and avoids asserting the identity for arbitrary accepted variables.
\end{proof}

\begin{remark}[Role relative to the Le Cam comparison]
\label{rem:accepted_transcript_vs_lecam_v10}
The lower bound in \cref{thm:two_point_lower_bound} is a deliberately restricted two-point benchmark.  \Cref{prop:accepted_transcript_efficiency_v10} is a sharper statement for the transcript actually used by SV in the bias-negligible regime.  Both are useful, but they answer different questions: Le Cam tests whether any estimator can distinguish two nearby ideal targets from noisy data, while accepted-transcript efficiency tests whether a chosen postselected estimator wastes information inside its declared transcript.
\end{remark}

\section{Toy model with closed constants and finite-\texorpdfstring{$p$}{p} boundary}
\label{sec:toy_model_v4}

This section gives a minimal analytic example.  It is not intended as a realistic benchmark; it demonstrates that the constants and validity certificates are computable on a finite interval, not only as formal Taylor coefficients.

Let the ideal one-qubit state be \(\rho_\star=\ket0\bra0\), the observable be \(O=X\), so \(\mu_\star=0\), and let the noisy state have spectral form
\begin{equation}
    \rho_p=(1-\varepsilon p)\ket{\phi_1(p)}\bra{\phi_1(p)}+
    \varepsilon p\ket{\phi_2(p)}\bra{\phi_2(p)},
    \label{eq:toy_state_v4}
\end{equation}
where
\begin{equation}
    \ket{\phi_1(p)}=\cos(\theta p)\ket0+\sin(\theta p)\ket1,
    \qquad
    \ket{\phi_2(p)}=-\sin(\theta p)\ket0+\cos(\theta p)\ket1.
\end{equation}
Assume \(0<p\leq p_{\max}\), where
\begin{equation}
    p_{\max}\leq \min\left\{\frac{1}{4\varepsilon},\frac{1}{4\abs\theta}\right\}.
    \label{eq:pmax_toy_v4}
\end{equation}
This keeps the spectrum ordered and the trigonometric expansions uniformly controlled.

\paragraph{Unmitigated estimator.}
Since \(\bra{\phi_1}X\ket{\phi_1}=\sin(2\theta p)\) and \(\bra{\phi_2}X\ket{\phi_2}=-\sin(2\theta p)\),
\begin{equation}
    \mu_0(p)=(1-2\varepsilon p)\sin(2\theta p).
\end{equation}
Using \(\abs{\sin x-x}\leq \abs{x}^3/6\), \(\abs{2\theta p}\le1/2\), and \(1-2\varepsilon p\ge1/2\), the unmitigated squared bias obeys the finite-interval bounds
\begin{equation}
    \frac12\theta^2p^2
    \leq
    b_0(p)^2
    \leq
    16\theta^2p^2
    \qquad (0<p\leq p_{\max}).
    \label{eq:toy_b0_bounds_v4}
\end{equation}
Indeed, the sharper lower constant \(529/576\) also follows from \(\abs{\sin x}\ge(23/24)\abs{x}\) on \(\abs{x}\le1/2\), but the displayed \(1/2\) is used below to keep the crossing certificate conservative.
The one-shot variance for measuring \(X\) is
\begin{equation}
    \sigma_0^2(p)=1-\mu_0(p)^2,
\end{equation}
and hence
\begin{equation}
    1-16\theta^2p_{\max}^2\leq \sigma_0^2(p)\leq1.
    \label{eq:toy_s0_bounds_v4}
\end{equation}

\paragraph{Virtual distillation.}
For \(M\geq2\),
\begin{equation}
    \mu_{\mathrm{VD},M}(p)=
    \frac{(1-\varepsilon p)^M-(\varepsilon p)^M}{(1-\varepsilon p)^M+(\varepsilon p)^M}\sin(2\theta p).
\end{equation}
The dominant-eigenvector floor is
\begin{equation}
    b_{\mathrm{VD},\infty}(p)=\sin(2\theta p)=2\theta p+\OO(p^3),
\end{equation}
so \(\delta_{\mathrm{VD}}=2\theta\).  VD cannot remove this coherent eigenvector mismatch.  The denominator is
\begin{equation}
    D_M(p)=(1-\varepsilon p)^M+(\varepsilon p)^M.
\end{equation}
For \(p\leq p_{\max}\),
\begin{equation}
    (1-\varepsilon p)^M\leq D_M(p)\leq (1-\varepsilon p)^M\left[1+\left(\frac{\varepsilon p}{1-\varepsilon p}\right)^M\right].
\end{equation}
The Hadamard denominator threshold \cref{eq:hadamard_Bmin_v4} is therefore explicit in \((\varepsilon,\theta,p,M)\).

\paragraph{Symmetry verification.}
Let \(P=\ket0\bra0\).  In this first toy, \(X\) is not a symmetry-compatible observable in the usual sense because \(XP\ne PX\).  The construction should therefore be read as an algebraic accepted-transcript quotient used to test the finite-shot crossing constants, not as a complete physical SV protocol.  The physically sector-compatible comparison is given in the next toy model.  Then
\begin{equation}
    a(p)=\Tr(P\rho_p)=
    (1-\varepsilon p)\cos^2(\theta p)+\varepsilon p\sin^2(\theta p).
\end{equation}
For \(p\leq p_{\max}\),
\begin{equation}
    \frac{11}{16}\leq a(p)\leq1.
    \label{eq:toy_acceptance_bounds_v4}
\end{equation}
Indeed, \(1-a(p)=\sin^2(\theta p)+\varepsilon p\cos(2\theta p)\le \theta^2p^2+\varepsilon p\le5/16\).
Moreover,
\begin{equation}
    \mu_{\mathrm{SV}}(p)=\frac{\Tr(XP\rho_pP)}{a(p)}=0,
\end{equation}
so \(\delta_{\mathrm{SV}}=0\).  If the accepted conditional \(X\)-measurement variance is one, then
\begin{equation}
    s_{\mathrm{SV}}(p)=\frac1{a(p)},
    \qquad
    1\leq s_{\mathrm{SV}}(p)\leq2.
    \label{eq:toy_sv_variance_bounds_v4}
\end{equation}


\paragraph{Toy remainder envelope.}
For this first toy model, the unmitigated estimator is the ordinary sample mean of a \(\{\pm1\}\)-valued variable.  Hence its expectation and variance are exactly
\begin{equation}
    \E\widehat\mu_0=\mu_0,
    \qquad
    \Var(\widehat\mu_0)=\sigma_0^2/B,
\end{equation}
so the unmitigated \(R^{-2}\) coefficient is
\begin{equation}
    C_{0,\mathrm{toy}}=0.
\end{equation}
For SV, the worst-case constants obtained from the general quotient certificate are deliberately loose.  In this first toy one can compute the relevant second-order remainder directly.  Let
\begin{equation}
    Z_s=A_s X_s,
    \qquad
    A_s\in\{0,1\},
    \qquad
    \widehat\mu_{\mathrm{SV}}=\frac{\overline Z}{\overline A},
\end{equation}
where \(A_s\) is the symmetry-acceptance indicator and \(X_s\in\{\pm1\}\) is the accepted \(X\)-measurement sign.  In this toy \(\mu_{\mathrm{SV}}=0\), \(\E Z_s=0\), \(\E A_s=a\), and
\begin{equation}
    \sigma_Z^2=a,
    \qquad
    \sigma_A^2=a(1-a),
    \qquad
    \sigma_{ZA}=0.
\end{equation}
Writing \(\bar z=\overline Z\) and \(\bar\alpha=\overline A-a\), the Neumann expansion gives
\begin{equation}
    \frac{\overline Z}{\overline A}
    =L_{\mathrm{SV}}+Q_{\mathrm{SV}}+\mathrm{higher\ order},
    \qquad
    L_{\mathrm{SV}}=\frac{\bar z}{a},
    \qquad
    Q_{\mathrm{SV}}=-\frac{\bar z\bar\alpha}{a^2}.
\end{equation}
A direct calculation of the sample-average moments gives the asymptotic second-order coefficient
\begin{equation}
    B^2\Var(Q_{\mathrm{SV}})=\frac{1-a}{a^2}+\OO(B^{-1}),
    \qquad
    2B^2\Cov(L_{\mathrm{SV}},Q_{\mathrm{SV}})=-\frac{2(1-a)}{a^2}.
    \label{eq:toy_sv_second_order_exact_final}
\end{equation}
Thus the second-order asymptotic variance coefficient is
\begin{equation}
    W_{\mathrm{SV}}=-\frac{1-a}{a^2},
    \qquad
    |W_{\mathrm{SV}}|\le2
    \quad\text{whenever }a\ge\frac12.
    \label{eq:toy_sv_W_bound_final}
\end{equation}
Because \(Z\in\{-1,0,1\}\), \(A\in\{0,1\}\), and \(a\ge1/2\), the cubic and higher Neumann terms can be bounded directly rather than by the general worst-case SV quotient certificate.

\begin{lemma}[Direct SV toy remainder envelope]
\label{lem:sv_toy_direct_remainder_v13}
In the first toy model, assume \(a\ge1/2\) and let \(R=B\) be the accepted-transcript shot budget.  On the good event
\begin{equation}
    \mathcal G=\{\abs{\overline A-a}\le a/2\},
\end{equation}
the Neumann remainder after the linear and quadratic terms satisfies
\begin{equation}
    \E\!\left[\1_{\mathcal G}\left(\frac{\overline Z}{\overline A}-L_{\mathrm{SV}}-Q_{\mathrm{SV}}\right)^2\right]
    \le \frac{8}{R^2}.
    \label{eq:sv_toy_good_remainder_v13}
\end{equation}
Moreover, by the acceptance Bernstein bound,
\begin{equation}
    \Pr(\mathcal G^c)\le 2e^{-R/16},
\end{equation}
and the clipped quotient is bounded by \(2\) on \(\mathcal G^c\).  Hence, once \(R\ge R_{\min,\mathrm{toy}}\) is chosen so that \(16e^{-R/16}\le 2/R^2\), the total additional contribution of the higher Neumann terms and the bad-denominator event is bounded by \(10/R^2\).
\end{lemma}

\begin{proof}
On \(\mathcal G\),
\begin{equation}
    \frac{\overline Z}{\overline A}
    =\frac{\bar z}{a}\sum_{k=0}^\infty\left(-\frac{\bar\alpha}{a}\right)^k,
    \qquad
    \abs{\bar\alpha/a}\le\frac12.
\end{equation}
Therefore the post-quadratic remainder is
\begin{equation}
    H_{\mathrm{SV}}
    :=\frac{\overline Z}{\overline A}-L_{\mathrm{SV}}-Q_{\mathrm{SV}}
    =\frac{\bar z\bar\alpha^2}{a^3}\frac{1}{1+\bar\alpha/a},
\end{equation}
and \(\abs{H_{\mathrm{SV}}}\le16\abs{\bar z}\bar\alpha^2\).  We now spell out the only contraction structure used in the displayed bound.  Since \(Z_sA_s=Z_s\), \(Z_s=0\) on rejection, and \(\E Z_s=0\), one has
\begin{equation}
    \E[Z_s(A_s-a)^r]=(1-a)^r\E Z_s=0,
    \qquad r=0,1,2,3,4.
    \label{eq:single_Z_block_v14}
\end{equation}
Thus any sample-index block containing a single \(Z\) and any number of \(\alpha=A-a\) factors vanishes.  In the expansion of \(\bar z^{\,2}\bar\alpha^4\), the two \(Z\)-indices must therefore coincide; after this identification, every remaining sample index carrying only \(\alpha\)-factors must occur with multiplicity at least two, because \(\E(A_s-a)=0\).  The surviving partitions are consequently only \(\{ZZ\}\{\alpha\alpha\}\{\alpha\alpha\}\), \(\{ZZ\alpha\}\{\alpha\alpha\alpha\}\), \(\{ZZ\alpha\alpha\}\{\alpha\alpha\}\), and \(\{ZZ\alpha\alpha\alpha\alpha\}\), together with index coincidences among these blocks.  Using \(Z_s^2\le A_s\), \(\abs{A_s-a}\le1\), and \(a\ge1/2\), this finite enumeration gives
\begin{equation}
    \E[\bar z^{\,2}\bar\alpha^{4}]
    \le \frac{1}{32R^2}.
    \label{eq:sv_toy_contraction_bound_v13}
\end{equation}
Combining the previous two displays gives \cref{eq:sv_toy_good_remainder_v13}.  Finally, on \(\mathcal G^c\) the clipped quotient contribution is bounded by \(4\Pr(\mathcal G^c)\le8e^{-R/16}\); the stated slightly looser absorption \(16e^{-R/16}\le2/R^2\) leaves room for the corresponding expectation/variance cross terms.  Thus \(8/R^2+2/R^2=10/R^2\), proving the claimed envelope.
\end{proof}

Consequently, for the certified toy crossing on budgets \(R\ge R_{\min,\mathrm{toy}}\), an admissible direct remainder envelope is
\begin{equation}
    C_{\mathrm{SV},\mathrm{toy}}^{\mathrm{op}}:=10.
    \label{eq:Ctoy_operational_final}
\end{equation}
This replaces the earlier worst-case envelope of order \(10^7\) in this specific toy.  The large general constant remains a valid universal certificate, but it is not the operational constant governing this closed model.

\begin{proposition}[Two-sided finite-\(p\) SV/no-mitigation boundary]
\label{prop:toy_two_sided_boundary_v4}
For the toy model above, assume \(p\in(0,p_{\max}]\), \(16\theta^2p_{\max}^2\leq1/2\), resources are restricted to \(R\ge R_{\min,\mathrm{toy}}\), and resource normalization \(\kappa_0=\kappa_{\mathrm{SV}}=1\).  The leading SV/no-mitigation crossing
\begin{equation}
    R_{\mathrm{SV}\leftrightarrow0}(p)
    =\frac{s_{\mathrm{SV}}(p)-s_0(p)}{b_0(p)^2-b_{\mathrm{SV}}(p)^2}
\end{equation}
satisfies the finite-interval sandwich
\begin{equation}
    \frac{\varepsilon}{64\theta^2}\frac1p
    \leq
    R_{\mathrm{SV}\leftrightarrow0}(p)
    \leq
    \frac{8\varepsilon+64\theta^2p_{\max}}{\theta^2}\frac1p.
    \label{eq:toy_boundary_sandwich_v4}
\end{equation}
In particular the boundary scales as \(1/p\) uniformly on the declared interval.
\end{proposition}

\begin{proof}
The denominator of the crossing is \(b_0(p)^2\), because \(b_{\mathrm{SV}}=0\).  The two-sided bias control is \cref{eq:toy_b0_bounds_v4}.  For the numerator,
\begin{equation}
    s_{\mathrm{SV}}(p)-s_0(p)
    =\frac1{a(p)}-\bigl(1-\mu_0(p)^2\bigr)
    =\frac{1-a(p)}{a(p)}+\mu_0(p)^2.
\end{equation}
Furthermore,
\begin{equation}
    1-a(p)
    =\sin^2(\theta p)+\varepsilon p\cos(2\theta p).
\end{equation}
Since \(p\leq1/(4|\theta|)\), \(\cos(2\theta p)\geq1/2\).  Hence
\begin{equation}
    s_{\mathrm{SV}}(p)-s_0(p)\geq 1-a(p)\geq \frac{\varepsilon p}{2}.
    \label{eq:toy_numerator_lower_fixed}
\end{equation}
For the upper bound, \(a(p)\geq1/2\), \(\mu_0(p)^2\leq16\theta^2p^2\), and
\[
    1-a(p)\leq \varepsilon p+\theta^2p^2
\]
give the conservative estimate
\begin{equation}
    s_{\mathrm{SV}}(p)-s_0(p)
    \leq 4\varepsilon p+32\theta^2p_{\max}p.
    \label{eq:toy_numerator_upper_fixed}
\end{equation}
The constants in the last display are intentionally loose but uniform on the declared interval.  Dividing \cref{eq:toy_numerator_lower_fixed,eq:toy_numerator_upper_fixed} by the upper and lower bounds on \(b_0(p)^2\), respectively, and adding one harmless factor-two slack to the lower side, gives \cref{eq:toy_boundary_sandwich_v4}.  In particular, the \(p\)-linear acceptance penalty divided by the \(p^2\)-bias gap gives the \(1/p\) boundary.
\end{proof}


\begin{proposition}[Operationally non-empty certified regime for the first toy crossing]
\label{prop:toy_certified_nonempty_v5}
Under the assumptions of \cref{prop:toy_two_sided_boundary_v4}, use the direct SV toy remainder envelope \(C_{\mathrm{SV},\mathrm{toy}}^{\mathrm{op}}=10\) from \cref{eq:Ctoy_operational_final}.  The certified crossing theorem applies whenever
\begin{equation}
    \eta_{\mathrm{toy}}
    :=\frac{4C_{\mathrm{SV},\mathrm{toy}}^{\mathrm{op}}}{g(p)R_{\mathrm{SV}\leftrightarrow0}(p)^2}<\frac12,
    \qquad
    g(p)=b_0(p)^2.
    \label{eq:toy_eta_verified_condition_v5}
\end{equation}
A sufficient parameter-level condition, using the finite-interval sandwich of \cref{prop:toy_two_sided_boundary_v4}, is
\begin{equation}
    \frac{32768\,C_{\mathrm{SV},\mathrm{toy}}^{\mathrm{op}}\theta^2}{\varepsilon^2}<\frac12.
    \label{eq:toy_explicit_nonempty_condition_v5}
\end{equation}
For example,
\begin{equation}
    \varepsilon=1,
    \qquad
    |\theta|=10^{-3},
    \qquad
    p_{\max}=10^{-1}
\end{equation}
satisfy the small-angle assumptions and the certified-crossing condition, because
\begin{equation}
    32768\cdot 10\cdot10^{-6}=3.2768\times10^{-1}<\frac12.
\end{equation}
At \(p=p_{\max}\), the lower sandwich gives \(R_{\mathrm{SV}\leftrightarrow0}\ge1/(64\cdot10^{-6}\cdot10^{-1})>1.5\times10^5\), which is safely above the fixed exponential-absorption threshold \(R_{\min,\mathrm{toy}}\).  Thus the certified region is non-empty at ordinary small-noise parameters.  The improvement over the earlier \(\theta\sim10^{-6}\) check comes from using the direct second-order quotient remainder in this closed toy, not from the Efron--Stein total-variance bound.
\end{proposition}

\begin{proof}
For this toy crossing, \(g(p)=b_0(p)^2-b_{\mathrm{SV}}(p)^2=b_0(p)^2\) and \(R_0=R_{\mathrm{SV}\leftrightarrow0}(p)\).  The certified-crossing parameter in \cref{thm:certified_crossing} is \(4C/[g(p)R_0^2]\).  The sandwich bounds give \(g(p)\ge\theta^2p^2/2\) and
\begin{equation}
    R_0\ge\frac{\varepsilon}{64\theta^2p}.
\end{equation}
Therefore
\begin{equation}
    g(p)R_0^2
    \ge
    \frac{\theta^2p^2}{2}\frac{\varepsilon^2}{4096\theta^4p^2}
    =\frac{\varepsilon^2}{8192\theta^2},
\end{equation}
and hence
\begin{equation}
    \eta_{\mathrm{toy}}
    \le
    \frac{4C\cdot8192\theta^2}{\varepsilon^2}
    =\frac{32768C\theta^2}{\varepsilon^2}.
\end{equation}
Substituting \(C=C_{\mathrm{SV},\mathrm{toy}}^{\mathrm{op}}=10\), \(\varepsilon=1\), and \(|\theta|=10^{-3}\) proves the displayed numerical certificate.  Finally, \(p_{\max}=10^{-1}\) obeys \(p_{\max}\le1/(4\varepsilon)\) and \(16\theta^2p_{\max}^2\ll1/2\), so the assumptions of the finite-interval sandwich are satisfied.
\end{proof}



\begin{proposition}[Selector versus Le Cam/Fisher benchmarks in the one-qubit transcript]
\label{prop:selector_lecam_constant_factor_v8}
\textup{(Conservative and sharpened versions.)}
View the first toy model through the same one-qubit accepted-\(X\)-measurement transcript used in \cref{prop:explicit_ip_bernoulli_v5}.  Assume the certified crossing regime of \cref{prop:toy_certified_nonempty_v5}, and restrict to the SV side of that certified crossing, i.e.\ budgets \(R\) larger than the certified SV/no-mitigation crossing.  Then the selected leading law is
\begin{equation}
    \MSE_{\rm sel}^{\rm lead}(p,R)=\frac{s_{\mathrm{SV}}(p)}{R}=\frac{1}{a(p)R}.
\end{equation}
Against the corrected conservative finite-interval Le Cam lower bound
\begin{equation}
    L_{\rm LC}^{\rm cons}(p,R)=\frac{1}{32R(1-p)^2},
\end{equation}
one has
\begin{equation}
    \frac{\MSE_{\rm sel}^{\rm lead}(p,R)}{L_{\rm LC}^{\rm cons}(p,R)}
    =\frac{32(1-p)^2}{a(p)}\le64.
    \label{eq:selector_lecam_ratio_v7}
\end{equation}
Using instead the local Fisher/two-point constant \(I_{p,\mathrm{loc}}^{\rm Bern}=2(1-p)^2\) from \cref{eq:explicit_ip_fisher_local_v8}, the corresponding local Le Cam scale is
\begin{equation}
    L_{\rm LC}^{\rm loc}(p,R)=\frac{1}{16R(1-p)^2},
\end{equation}
and therefore
\begin{equation}
    \frac{\MSE_{\rm sel}^{\rm lead}(p,R)}{L_{\rm LC}^{\rm loc}(p,R)}
    =\frac{16(1-p)^2}{a(p)}\le32.
    \label{eq:selector_lecam_ratio_local_v8}
\end{equation}
Finally, if the same accepted transcript is benchmarked by its regular local asymptotic Fisher information \(J_{\rm acc}(p)=a(p)\) for the accepted ideal \(X\)-mean in this deliberately unbiased SV toy, then the Cramer--Rao/LAN scale is
\begin{equation}
    L_{\rm LAN}^{\rm acc}(p,R)=\frac{1}{a(p)R},
\end{equation}
and the selected SV law saturates this accepted-transcript Fisher scale exactly.  Thus the factors \(64\) and \(32\) in \cref{eq:selector_lecam_ratio_v7,eq:selector_lecam_ratio_local_v8} are conservative two-point artifacts, not physical evidence that the selector is far from the relevant sampling scale.
\end{proposition}

\begin{proof}
In the first toy, \(\mu_{\mathrm{SV}}(p)=\mu_\star=0\), so the leading SV MSE is purely sampling and equals \(s_{\mathrm{SV}}(p)/R=1/[a(p)R]\).  The phrase ``selected'' is justified only on the SV side of the certified crossing supplied by \cref{prop:toy_certified_nonempty_v5}; outside that regime the statement is only a comparison of the SV estimator, not of the selector.  Dividing \(1/[a(p)R]\) by the corrected conservative and local Le Cam scales gives \cref{eq:selector_lecam_ratio_v7,eq:selector_lecam_ratio_local_v8}; \(a(p)\ge1/2\) follows from \cref{eq:toy_acceptance_bounds_v4}.  The number of accepted samples concentrates around \(a(p)R\) by the same multiplicative Chernoff bound used in \cref{thm:sv_quotient_law_v2}; hence the phrase \(a(p)R\) effective observations can be made high-probability rather than merely heuristic.  The final LAN statement is the ordinary Fisher benchmark for those accepted unit-variance observations.  It is included only to show what part of the constant factor comes from deliberately conservative two-point bounds.
\end{proof}

\begin{remark}[Role of the first toy model]
The first toy model is deliberately favorable to SV: \(\delta_{\mathrm{SV}}=0\) because the rank-one projector kills the \(X\)-expectation inside the accepted block.  Its purpose is not to prove a generic VD--SV comparison; its purpose is to show that the finite-shot crossing certificate can be checked on a finite interval with explicit constants.  The next toy model gives the promised nontrivial VD--SV comparison with \(\delta_{\mathrm{SV}}\ne0\).
\end{remark}



\subsection{A sector toy model with \texorpdfstring{$\delta_{\mathrm{SV}}\ne0$}{deltaSV nonzero}}
\label{subsec:sector_toy_delta_sv_nonzero_v5}
The previous example has a rank-one symmetry sector, so SV has zero residual first-order bias.  The present sector model is the complementary clean limit: it has nonzero SV residual bias but no dominant-eigenvector rotation, so \(\delta_{\mathrm{VD}}=0\) at first order.  These examples are deliberately diagnostic, not generic; a fully generic toy model can have both \(\delta_{\mathrm{VD}}\ne0\) and \(\delta_{\mathrm{SV}}\ne0\).  To exhibit a genuine VD--SV boundary, consider a three-dimensional Hilbert space with orthonormal vectors \(\ket g,\ket e,\ket\ell\).  The valid sector is \(P=\ket g\bra g+\ket e\bra e\), the ideal state is \(\rho_\star=\ket g\bra g\), and the observable is
\begin{equation}
    O=\ket g\bra g-\ket e\bra e,
    \qquad O\ket\ell=0,
    \qquad \mu_\star=1.
\end{equation}
Let
\begin{equation}
    \rho_p=\alpha(p)\ket g\bra g+\beta(p)\ket e\bra e+\lambda(p)\ket\ell\bra\ell,
    \label{eq:sector_toy_state_v5}
\end{equation}
with
\begin{equation}
    \alpha(p)=1-(u+\ell)p,
    \qquad
    \beta(p)=up,
    \qquad
    \lambda(p)=\ell p,
\end{equation}
where \(u,\ell>0\) and \(0<p\le p_{\max}\le[4(u+\ell)]^{-1}\).  Here \(up\) is undetectable in-sector error and \(\ell p\) is detectable leakage.

\paragraph{SV constants.}
The acceptance probability and SV population value are
\begin{equation}
    a(p)=\alpha(p)+\beta(p)=1-\ell p,
    \qquad
    \mu_{\mathrm{SV}}(p)=\frac{\alpha(p)-\beta(p)}{\alpha(p)+\beta(p)}
    =1-\frac{2up}{1-\ell p}.
\end{equation}
Therefore
\begin{equation}
    b_{\mathrm{SV}}(p)=-\frac{2up}{1-\ell p},
    \qquad
    \delta_{\mathrm{SV}}=-2u\ne0.
    \label{eq:sector_toy_delta_sv_v5}
\end{equation}
The accepted-shot variance coefficient is
\begin{equation}
    v_{\mathrm{SV}}(p)=\frac{1-\mu_{\mathrm{SV}}(p)^2}{a(p)}
    =\frac{4up(1-(u+\ell)p)}{(1-\ell p)^3}
    =4up+\OO(p^2).
    \label{eq:sector_toy_vsv_v5}
\end{equation}

\paragraph{VD constants.}
For VD with copy number \(M\),
\begin{equation}
    D_M(p)=\alpha(p)^M+\beta(p)^M+\lambda(p)^M,
    \qquad
    N_M(p)=\alpha(p)^M-\beta(p)^M,
\end{equation}
so
\begin{equation}
    \mu_{\mathrm{VD},M}(p)=\frac{\alpha(p)^M-\beta(p)^M}{\alpha(p)^M+\beta(p)^M+\lambda(p)^M}.
\end{equation}
For \(M=2\),
\begin{align}
    D_2(p)&=\alpha^2+\beta^2+\lambda^2
    =1-2(u+\ell)p+\OO(p^2),\label{eq:sector_D2_expansion_v6}\\
    N_2(p)&=\alpha^2-\beta^2
    =1-2(u+\ell)p+\OO(p^2).\label{eq:sector_N2_expansion_v6}
\end{align}
A direct subtraction gives
\begin{equation}
    b_{\mathrm{VD},2}(p)
    =\frac{N_2(p)}{D_2(p)}-1
    =-\frac{2\beta(p)^2+\lambda(p)^2}{D_2(p)}
    =-\frac{(2u^2+\ell^2)p^2}{D_2(p)}
    =-(2u^2+\ell^2)p^2+\OO(p^3),
    \label{eq:sector_toy_bvd2_v5}
\end{equation}
so VD removes the first-order in-sector error by spectral amplification.  Under the independent Hadamard-test implementation of \cref{prop:implementation_variance_inflation_v4},
\begin{equation}
    v_{\mathrm{VD},2}(p)
    =\frac{1-N_2(p)^2+\mu_{\mathrm{VD},2}(p)^2(1-D_2(p)^2)}{D_2(p)^2},
    \label{eq:sector_toy_vvd2_formula_v6}
\end{equation}
where the covariance term is zero because the numerator and denominator Hadamard tests are independent.  Since \(\mu_{\mathrm{VD},2}(p)=1+\OO(p^2)\), \cref{eq:sector_D2_expansion_v6,eq:sector_N2_expansion_v6} imply
\begin{equation}
    1-D_2(p)^2=4(u+\ell)p+\OO(p^2),
    \qquad
    1-N_2(p)^2=4(u+\ell)p+\OO(p^2),
\end{equation}
and therefore
\begin{equation}
    v_{\mathrm{VD},2}(p)
    =8(u+\ell)p+\OO(p^2).
    \label{eq:sector_toy_vvd2_v5}
\end{equation}


\paragraph{Sector remainder and denominator certificates.}
On \(0<p\le p_{\max}\le[4(u+\ell)]^{-1}\), one has \(\alpha(p)\ge3/4\) and hence
\begin{equation}
    D_2(p)=\alpha^2+\beta^2+\lambda^2\ge \frac{9}{16},
    \qquad
    a(p)=1-\ell p\ge\frac34.
    \label{eq:sector_D_a_lower_v6}
\end{equation}
For the independent Hadamard-test VD implementation with \(M=2\), take \(K_{N,2}=1\), \(K_{D,2}=2\), and the bounded-moment envelopes
\begin{equation}
    \mathfrak m_{12},\mathfrak m_{21},\mathfrak m_{03}\le 8,
    \qquad
    \mathfrak m_{13},\mathfrak m_{04},\mathfrak m_{22},\mathfrak m_{40}\le 64,
    \qquad
    \mathfrak t^{(V)}\le 10^6.
\end{equation}
The final bound on \(\mathfrak t^{(V)}\) is a direct bounded-sign envelope for the third-and-higher quotient terms in this finite sector toy; it is deliberately loose and remains far below the already conservative crossing constant used below.
The finite-interval constant is assembled explicitly as follows.  Let
\(\Delta=(9/16)^{-1}=16/9\).  Since \(D_2^{-1}\le\Delta\) and
\(|\mu_{\mathrm{VD},2}|\le1\), \cref{eq:AE_constant_v4} gives
\begin{align}
    A_{E,\mathrm{VD},2}^{\mathrm{sector}}
    &\le (8+8)\Delta^3+64\Delta^4+2\cdot1\cdot64\Delta^5 \\
    &<3.1\times10^3.
    \label{eq:sector_AE_arithmetic_v7}
\end{align}
Using the same envelopes in \cref{eq:AQ_constant_v14,eq:ALQ_constant_v14,eq:AV_constant_v4},
\begin{align}
    A_{V,\mathrm{VD},2}^{\mathrm{sector}}
    &\le 2\Delta^4(64+64)
      +2\Delta^3(2\cdot8+8+8)+10^6\\
    &<1.1\times10^6<2.5\times10^8.
    \label{eq:sector_AV_arithmetic_v7}
\end{align}
This is the advertised conservative finite-interval bound
\begin{equation}
    C_{\mathrm{VD},2}^{\mathrm{sector}}\le 2.5\times10^8.
    \label{eq:Cvd_sector_v6}
\end{equation}
For SV in the sector toy, \(a\ge3/4\), \(K_W=1\), and the Bernoulli quotient constants give
\begin{equation}
    C_{\mathrm{SV}}^{\mathrm{sector}}\le 5.0\times10^7.
    \label{eq:Csv_sector_v6}
\end{equation}
We therefore use the explicit envelope
\begin{equation}
    C_{\mathrm{sector}}:=3.0\times10^8.
    \label{eq:Csector_numeric_v6}
\end{equation}
The internal VD denominator certificate is also explicit.  Since \(\sigma_{D,2}^2\le1\), \(K_{D,2}=2\), and \(D_2\ge9/16\), \cref{eq:denom_sample_threshold_v4} gives
\begin{equation}
    B_{\mathrm{den},2}(p,\varepsilon_{\rm fail})
    \le 31\log\frac1{\varepsilon_{\rm fail}}.
    \label{eq:sector_Bden_bound_v6}
\end{equation}
Thus the VD law used inside the crossing is certified whenever the VD sample budget at the predicted crossing is at least the right-hand side of \cref{eq:sector_Bden_bound_v6}.

\begin{proposition}[Nontrivial VD--SV crossing in the sector toy]
\label{prop:sector_toy_vd_sv_crossing_v5}
In the sector toy, with \(M=2\), the leading VD--SV crossing satisfies
\begin{equation}
    R_{\mathrm{VD},2\leftrightarrow\mathrm{SV}}(p)
    =\frac{v_{\mathrm{VD},2}(p)-v_{\mathrm{SV}}(p)}{b_{\mathrm{SV}}(p)^2-b_{\mathrm{VD},2}(p)^2}
    =\frac{u+2\ell}{u^2}\frac1p+\OO(1).
    \label{eq:sector_toy_vd_sv_crossing_v5}
\end{equation}
Moreover, using the explicit envelope \(C_{\mathrm{sector}}=3.0\times10^8\), the crossing is certified whenever
\begin{equation}
    \eta_{\mathrm{sector}}
    =\frac{4C_{\mathrm{sector}}}{(b_{\mathrm{SV}}^2-b_{\mathrm{VD},2}^2)R_{\mathrm{VD},2\leftrightarrow\mathrm{SV}}(p)^2}<\frac12.
\end{equation}
Using the first-order bounds above, a sufficient non-empty condition is
\begin{equation}
    \frac{8C_{\mathrm{sector}}u^2}{(u+2\ell)^2}<\frac12.
    \label{eq:sector_toy_certificate_condition_v5}
\end{equation}
The internal VD quotient certificate is simultaneously valid if, with the chosen resource normalization,
\begin{equation}
    \frac{u+2\ell}{u^2p_{\max}}
    \ge 31\log\frac1{\varepsilon_{\rm fail}}.
    \label{eq:sector_internal_vd_certificate_v6}
\end{equation}
For example, \(u=1\), \(\ell=10^6\), \(p_{\max}=10^{-8}\), and \(\varepsilon_{\rm fail}=10^{-6}\) satisfy both \cref{eq:sector_toy_certificate_condition_v5,eq:sector_internal_vd_certificate_v6}.  This numerical point is deliberately extreme: \(\ell/u=10^6\) and \(p_{\max}=10^{-8}\) should not be read as a realistic device regime.  They only demonstrate that the nested certificates are jointly non-empty under the intentionally loose constants of \cref{rem:constant_tightness_v7}.  Sharper variance remainders would move the certified example toward much less extreme parameter ratios.  Hence the VD--SV crossing certificate and the internal VD denominator certificate are jointly non-empty.
\end{proposition}

\begin{proof}
Equations \cref{eq:sector_toy_delta_sv_v5,eq:sector_toy_bvd2_v5} give
\(b_{\mathrm{SV}}(p)^2-b_{\mathrm{VD},2}(p)^2=4u^2p^2+\OO(p^3)\).  Equations \cref{eq:sector_toy_vsv_v5,eq:sector_toy_vvd2_v5} give
\(v_{\mathrm{VD},2}(p)-v_{\mathrm{SV}}(p)=(4u+8\ell)p+\OO(p^2)\).  Dividing proves \cref{eq:sector_toy_vd_sv_crossing_v5}.  The certified condition is exactly \cref{eq:eta_crossing_condition} with the explicit sector-toy constants from \cref{eq:Csector_numeric_v6}.  Substituting the leading lower bound \(b_{\mathrm{SV}}^2-b_{\mathrm{VD},2}^2\ge 2u^2p^2\) and the crossing lower bound \(R_{\mathrm{VD},2\leftrightarrow\mathrm{SV}}\ge (u+2\ell)/(2u^2p)\) for small enough \(p_{\max}\) yields \cref{eq:sector_toy_certificate_condition_v5}.  The separate lower bound \cref{eq:sector_internal_vd_certificate_v6} is just \(R_{\mathrm{cross}}(p_{\max})\ge B_{\mathrm{den},2}\), using \cref{eq:sector_Bden_bound_v6}.  The displayed numerical choice verifies both inequalities directly.
\end{proof}


\begin{proposition}[Efron--Stein total-MSE dominance check at physical noise]
\label{prop:sector_es_operational_certificate_v12}
The extreme numerical point in \cref{prop:sector_toy_vd_sv_crossing_v5} is not needed to obtain an operationally meaningful dominance certificate.  In the sector toy with
\begin{equation}
    u=1,\qquad \ell=1,\qquad p=10^{-2},\qquad M=2,
\end{equation}
and unit Hadamard-test call costs, let \(R\) denote the physical call budget.  Since the independent VD implementation uses one numerator and one denominator stream, its paired-shot budget is \(B=R/2\).  The clipped VD estimator is certified to beat SV at any physical budget
\begin{equation}
    R\ge 2\times10^6.
    \label{eq:sector_es_budget_v12}
\end{equation}
More precisely, the Efron--Stein total-variance certificate gives
\begin{equation}
    \Var(\widehat\mu_{\mathrm{VD},2}^{\rm clip})
    \le \frac{472}{R},
    \label{eq:sector_es_236_v12}
\end{equation}
while the VD population bias satisfies
\begin{equation}
    \abs{b_{\mathrm{VD},2}(p)}\le 4p^2,
    \label{eq:sector_es_bvd_v12}
\end{equation}
and the SV residual floor satisfies
\begin{equation}
    \MSE_{\mathrm{SV}}(p,R)\ge b_{\mathrm{SV}}(p)^2\ge 4p^2.
    \label{eq:sector_es_sv_floor_v12}
\end{equation}
Consequently, at \(p=10^{-2}\) and \(R=2\times10^6\),
\begin{equation}
    \MSE_{\mathrm{VD},2}^{\rm cert}(p,R)
    \le (4p^2+4/R)^2+\frac{472}{R}+10^{-8}
    <2.4\times10^{-4}
    <4.0\times10^{-4}
    \le \MSE_{\mathrm{SV}}(p,R).
    \label{eq:sector_es_numeric_v12}
\end{equation}
Thus VD can be certified to beat SV in total MSE at a physical noise scale using modest sector ratios.  This is not a replacement for the crossing-remainder certificate in \cref{thm:certified_crossing}: Efron--Stein controls the total variance of the clipped quotient, not the difference between the exact MSE and its leading local law.  It is included only as a robust dominance check showing that the extreme constants in the formal crossing certificate are not physical thresholds.
\end{proposition}

\begin{proof}
For \(u=\ell=1\) and \(p=10^{-2}\),
\begin{equation}
    D_2=(1-2p)^2+p^2+p^2=0.9606>0.96.
\end{equation}
With \(K_{N,2}=1\) and \(K_{D,2}=2\), \cref{eq:efron_stein_vd_bound_v11} gives
\begin{equation}
    \frac{8K_{N,2}^2(D_2+2K_{D,2})^2}{D_2^4}
    \le \frac{8(1+4)^2}{0.96^4}<236,
\end{equation}
which gives \(\Var(\widehat\mu_{\mathrm{VD},2}^{\rm clip})\le236/B=472/R\) because \(B=R/2\).  This proves \cref{eq:sector_es_236_v12}.  The exact VD population bias in the sector toy is
\begin{equation}
    b_{\mathrm{VD},2}(p)=-\frac{(2u^2+\ell^2)p^2}{D_2(p)},
\end{equation}
so \(\abs{b_{\mathrm{VD},2}}\le 3p^2/0.96<4p^2\).  The SV bias is
\begin{equation}
    b_{\mathrm{SV}}(p)=-\frac{2up}{1-\ell p}.
\end{equation}
Here \(u=\ell=1\) and \(p=10^{-2}<1\), so \(0<1-\ell p<1\).  Hence
\begin{equation}
    b_{\mathrm{SV}}(p)^2
    =\frac{4p^2}{(1-p)^2}
    >4p^2.
    \label{eq:sector_sv_floor_margin_v14}
\end{equation}
Numerically this is \(4.0812\ldots\times10^{-4}\), safely above \(4.0\times10^{-4}\); the inequality follows from the denominator \(1-p<1\), not from rounding.  It remains to justify the displayed \(4/R\) finite-sample quotient-bias allowance.  In this independent Hadamard-test implementation \(\sigma_{ND}=0\) and the denominator sign has variance \(\sigma_D^2=1-D_2^2\).  Hence
\begin{equation}
    \frac{\abs{c_{\mathrm{VD},2}}}{B}
    =\frac{\abs{\mu_{\mathrm{VD},2}}\sigma_D^2}{D_2^2B}
    \le \frac{1-D_2^2}{D_2^2B}<\frac{0.09}{B}<\frac{4}{R}.
\end{equation}
The remaining expectation remainder is controlled by \cref{eq:sector_AE_arithmetic_v7}:
\begin{equation}
    \frac{A_{E,\mathrm{VD},2}^{\mathrm{sector}}}{B^2}
    <\frac{3.1\times10^3}{(10^6)^2}<3.1\times10^{-9}.
\end{equation}
For the denominator event, \cref{eq:sector_Bden_bound_v6} gives \(B_{\mathrm{den},2}\le31\log(10^6)<4.3\times10^2\) at failure probability \(10^{-6}\), while here \(B=10^6\); hence the exponentially small denominator contribution is also absorbed below \(10^{-8}\).  Substituting \(p=10^{-2}\) and \(R=2\times10^6\) yields \cref{eq:sector_es_numeric_v12}.
\end{proof}

\begin{proposition}[Sector residual-bias floor and a structural two-point benchmark]
\label{prop:sector_structural_floor_v8}
In the sector toy, SV has residual first-order bias
\begin{equation}
    b_{\mathrm{SV}}(p)=-2up+O(p^2),
    \qquad
    b_{\mathrm{SV}}(p)^2=4u^2p^2+O(p^3).
\end{equation}
Consider the restricted ambiguity model in which the same full noisy state \(\rho_p\) can be generated either from the ideal state \(\rho_+^\star=\ket g\bra g\) with an undetectable in-sector error of mass \(up\), or from an alternative ideal state \(\rho_-^\star=(1-up)\ket g\bra g+up\ket e\bra e\) with that in-sector mass treated as ideal population, with the remaining channel adjusted so that the transcript distribution is identical.  The two ideal values differ by
\begin{equation}
    \mu_+^\star-\mu_-^\star=2up.
\end{equation}
Hence \cref{cor:nonidentifiability_floor} gives the structural lower bound
\begin{equation}
    L_{\rm str}^{\rm sector}(p)=\frac{u^2p^2}{2}.
    \label{eq:sector_structural_floor_v8}
\end{equation}
Consequently the SV residual bias floor is within a constant factor of this restricted non-identifiability floor:
\begin{equation}
    \frac{b_{\mathrm{SV}}(p)^2}{L_{\rm str}^{\rm sector}(p)}
    =8+O(p).
    \label{eq:sector_sv_floor_ratio_v8}
\end{equation}
This does not say that VD cannot beat SV in the declared known-channel sector model; VD does beat SV on the high-budget side of \cref{prop:sector_toy_vd_sv_crossing_v5}.  It says something narrower: if the in-sector population error is not identifiable from side information, then the same order of bias floor is information-theoretic rather than merely a defect of SV\@.
\end{proposition}

\begin{proof}
The expansion of \(b_{\mathrm{SV}}\) is \cref{eq:sector_toy_delta_sv_v5}.  The two hypotheses described in the statement induce the same noisy transcript distribution by construction, while their ideal \(O=\ket g\bra g-\ket e\bra e\) values are \(1\) and \(1-2up\).  Applying \cref{cor:nonidentifiability_floor} with \(h=up\) gives \cref{eq:sector_structural_floor_v8}.  Dividing the SV squared bias by this floor gives \cref{eq:sector_sv_floor_ratio_v8}.
\end{proof}

\begin{remark}[Role of the sector toy model]
The sector toy is deliberately favorable to VD in a different way than the first toy is favorable to SV\@.  Because \(\rho_p\) remains diagonal in the fixed basis \(\{\ket g,\ket e,\ket\ell\}\), the dominant eigenvector does not rotate and \(\delta_{\mathrm{VD}}=0\).  Its purpose is to isolate the spectral-amplification advantage of VD against a nonzero SV residual bias.  The generic operating-window theory above allows both \(\delta_{\mathrm{VD}}\ne0\) and \(\delta_{\mathrm{SV}}\ne0\); these two closed models are diagnostic boundary cases used to check constants, not a claim that either favorable limit is typical.
\end{remark}


\subsection{A generic toy with simultaneous VD and SV first-order bias}
\label{subsec:generic_toy_v10}

The two closed toy models above deliberately isolate clean limiting mechanisms.  The first makes SV unbiased at first order; the sector model removes dominant-eigenvector rotation and therefore gives \(\delta_{\mathrm{VD}}=0\).  The following minimal model records the generic situation in which both first-order floors can be nonzero.  It is included as a structural diagnostic only.  The direct trace expansion and the formal VD--SV crossing are displayed below; only the sharp variance constants and a full certified-boundary check for this generic non-diagonal toy are left for a later version.

\begin{proposition}[A toy with simultaneous VD and SV first-order bias]
\label{prop:generic_toy_v10}
Work in \(\mathbb C^3\) with basis \(\{\ket0,\ket1,\ket2\}\), valid sector
\begin{equation}
    P=\ket0\bra0+\ket1\bra1,
\end{equation}
ideal state \(\rho_\star=\ket0\bra0\), and observable
\begin{equation}
    O=\ket0\bra0-\ket1\bra1+c\big(\ket0\bra1+\ket1\bra0\big),
    \qquad \mu_\star=1.
\end{equation}
Let
\begin{equation}
    \Delta
    =u(\ket1\bra1-\ket0\bra0)
    +\ell(\ket2\bra2-\ket0\bra0)
    +\kappa(\ket0\bra1+\ket1\bra0),
    \qquad u,\ell,\kappa>0 ,
\end{equation}
and set \(\rho_p=\rho_\star+p\Delta\) for
\begin{equation}
    0<p\leq p_0,
    \qquad
    p_0:=\frac12\min\left\{\frac1{u+\ell},\frac{u}{\kappa^2}\right\}.
    \label{eq:generic_toy_positivity_p0_v13}
\end{equation}
Then \(\rho_p\) is a density matrix on this interval: the \(\{\ket0,\ket1\}\) block has nonnegative determinant and the \(\ket2\) weight is \(\ell p\ge0\).  Here \(up\) is an undetectable in-sector population error, \(\ell p\) is detectable leakage, and \(\kappa p\) is an in-sector coherence that tilts the dominant eigenvector.  Then
\begin{equation}
    a(p)=1-\ell p+\OO(p^2),
    \qquad
    \delta_{\mathrm{SV}}=2(c\kappa-u),
    \qquad
    \delta_{\mathrm{VD}}=2c\kappa.
\end{equation}
Thus, whenever \(c\kappa\ne0\) and \(c\kappa\ne u\), both VD and SV have nonzero first-order residual bias.  For \(M=2\), direct trace expansion gives
\begin{align}
    D_2(p)&=\Tr(\rho_p^2)=1-2(u+\ell)p+\OO(p^2),
    \label{eq:generic_toy_D2_v12}\\
    N_2(p)&=\Tr(O\rho_p^2)=1+2(c\kappa-u-\ell)p+\OO(p^2),
    \label{eq:generic_toy_N2_v12}
\end{align}
and hence
\begin{equation}
    b_{\mathrm{VD},2}(p)=2c\kappa p+\OO(p^2),
    \qquad
    b_{\mathrm{SV}}(p)=2(c\kappa-u)p+\OO(p^2).
    \label{eq:generic_toy_biases_v12}
\end{equation}
Writing the leading variance coefficients as
\begin{equation}
    s_{\mathrm{VD},2}(p)=\nu_{\mathrm{VD},0}+\OO(p),
    \qquad
    s_{\mathrm{SV}}(p)=\nu_{\mathrm{SV},0}+\OO(p),
\end{equation}
the formal VD--SV crossing in the generic toy is
\begin{equation}
    R_{\mathrm{VD},2\leftrightarrow\mathrm{SV}}(p)
    =
    \frac{\nu_{\mathrm{VD},0}-\nu_{\mathrm{SV},0}+\OO(p)}
    {4[(c\kappa-u)^2-c^2\kappa^2]p^2+\OO(p^3)},
    \label{eq:generic_toy_crossing_v12}
\end{equation}
whenever the denominator is nonzero and the numerator has the sign required for a positive crossing.  Thus the generic simultaneous-bias case has the usual \(p^{-2}\) finite-shot scale, unlike the cleaner sector example where VD removes the first-order floor.
\end{proposition}

\begin{proof}
Since \(\Tr(P\Delta P)=\Delta_{00}+\Delta_{11}=-(u+\ell)+u=-\ell\), the acceptance probability is \(a(p)=1-\ell p+\OO(p^2)\).  Also
\begin{equation}
    \Tr(OP\Delta P)=O_{00}\Delta_{00}+O_{11}\Delta_{11}+2c\Delta_{01}=-(u+\ell)-u+2c\kappa.
\end{equation}
Using the first-order SV formula gives
\begin{equation}
    \delta_{\mathrm{SV}}
    =\Tr(OP\Delta P)-\mu_\star\Tr(P\Delta P)
    =-(2u+\ell)+2c\kappa+\ell
    =2(c\kappa-u).
\end{equation}
For VD, let \(Q_\star=I-\ket0\bra0\).  Then \(\Delta\ket0=-(u+\ell)\ket0+\kappa\ket1\), so \(Q_\star\Delta\ket0=\kappa\ket1\).  By \cref{prop:vd_delta_generator},
\begin{equation}
    \delta_{\mathrm{VD}}
    =2\Re\bra0 OQ_\star\Delta\ket0
    =2c\kappa.
\end{equation}
Finally, \(\rho_p=\rho_\star+p\Delta\) on the declared positivity interval gives
\begin{equation}
    \rho_p^2=\rho_\star+p(\rho_\star\Delta+\Delta\rho_\star)+\OO(p^2).
\end{equation}
Since \(\Tr(\rho_\star\Delta+\Delta\rho_\star)=2\Delta_{00}=-2(u+\ell)\), \cref{eq:generic_toy_D2_v12} follows.  Also
\begin{equation}
    \Tr\!\left[O(\rho_\star\Delta+\Delta\rho_\star)\right]
    =-2(u+\ell)+2c\kappa,
\end{equation}
which gives \cref{eq:generic_toy_N2_v12}.  Dividing \(N_2/D_2\) yields \(b_{\mathrm{VD},2}=2c\kappa p+\OO(p^2)\).  Combining this with the SV expansion above and equating the two leading MSE laws gives \cref{eq:generic_toy_crossing_v12}.
\end{proof}

\begin{remark}[Why the previous clean toys were special]
The first toy is SV-favourable because the postselected observable has zero residual first-order bias.  The sector toy is VD-favourable because the noisy state is diagonal in a fixed basis, so the dominant eigenvector does not rotate at first order.  \Cref{prop:generic_toy_v10} shows that neither simplification is generic: an in-sector coherence can make \(\delta_{\mathrm{VD}}\ne0\), and an undetectable in-sector population error can make \(\delta_{\mathrm{SV}}\ne0\).  In particular, a diagonal observable in the ideal eigenbasis has \(\delta_{\mathrm{VD}}=0\) because it cannot see the orthogonal first-order tilt direction \(Q_\star\Delta\ket{\psi_\star}\); off-diagonal observables generally can.
\end{remark}

\section{Reproducibility and data availability}

The accompanying artifact is organized so that the level of reproducibility matches the level of claim.  The Level~1 coefficient-checking layer is the main reproducible core: it uses a white-box density-matrix implementation, fixed instances and thresholds, and reference outputs checked by a verifier.  The fixed analysis plan has two archived hashes, and the amended v2 package passes a zero-drift check.  A reviewer can rerun the Level~1 reference verification without Qiskit or any quantum-cloud account.

Level~2 is a seeded Qiskit/Aer device-simulation layer.  Its outputs are included as a form-validation and robustness layer, but they are not mixed into the tight Level~1 reference tier.  They are documented as generated artifacts with seeds and reports, outside the Level~1 reference tier.

Level~3 is hardware analysis-reproducible only.  The raw IBM count archives, job identifiers, backend provenance, transpilation summaries, layouts, and calibration-epoch metadata are archived.  Public reproduction reruns only the analysis from those archives; it does not require, and should not attempt, new QPU submissions.  Hardware credentials are not stored in the repository.

The release package includes the code repository, archived analysis-plan hashes, reference outputs, archived hardware counts, exact dependency lockfiles, and release metadata.  The repository is public at \url{https://github.com/vicenzoscavino1999/finite-shot-vd-sv}, and the exact release used for this manuscript (\texttt{v1.0.0}) is archived at Zenodo under DOI \href{https://doi.org/10.5281/zenodo.20670996}{10.5281/zenodo.20670996}.

\section{Scope and limitations}

We prove conditional local laws, explicit quotient certificates, implementation-level VD variance constants for an idealized Hadamard/derangement estimator, a sharper Efron--Stein total-variance diagnostic, an asymptotic second-order quotient-variance diagnostic, finite-interval toy boundaries, a multiplicative implementation-noise baseline, and a restricted Le Cam lower bound.  The general toy certificates are numerically conservative; their large constants certify non-vacuity but are not intended as tight physical thresholds.  In the first closed toy, however, the direct SV second-order quotient calculation gives an operational remainder of order one and certifies the crossing at ordinary small-noise parameters.  The validation layer instantiates concrete QAOA instances, a concrete SV projector, realistic VD interferometry circuits, Aer device simulation, and qualitative IBM-hardware checks.  Those additions test form and feasibility, but they do not prove a global optimality theorem over all possible QEM protocols.  Such a theorem would require a fully specified statistical decision model for estimating an ideal functional \(\mu_\star\) from noisy states \(\rho_p\), including what is known about the noise channel and which adaptive measurements are admissible.

The remaining limitations are therefore methodological rather than missing implementation steps.
\begin{enumerate}[leftmargin=2em]
    \item The compiled QAOA/VD/SV circuits are used in the validation layer, but the theorem-level implementation-noise model remains the idealized multiplicative baseline.
    \item Level~2 and Level~3 use concrete QAOA families and readout-calibrated sector postselection, but they validate form and robustness rather than analytic coefficients.
    \item The hardware conclusion concerns the global interferometric VD circuits implemented here.  Localized virtual distillation, layout-aware compilation, or other lower-depth purification circuits can change implementation constants such as \(\kappa_M\) and move a crossing window; they do not remove the need to certify denominator concentration, quotient bias, and physical resource normalization.
    \item CDR remains outside the theorem-level selector unless a separate training-model law is proved.
    \item The hardware layer is qualitative and budget-limited: it uses three archived instances on a primary backend plus one independent cross-backend replicate of \texttt{g6\_k2\_c000}, with cluster bootstrap over calibration epochs, so it supports feasibility and robustness rather than a hardware-universal ordering.
\end{enumerate}

\section{Conclusion}

The strengthened mathematical core is this: VD and SV can be compared by finite-shot operating windows only after their residual bias floors, sampling coefficients, quotient corrections, and remainder regimes are all explicit.  In the closed one-qubit transcript, the selector can also be compared directly with conservative Le Cam and sharper local Fisher benchmarks; in the sector toy, the residual SV bias floor is compared with a restricted structural non-identifiability floor.  For VD, the key object is the quotient law plus the denominator certificate.  For SV, the key object is the projected first-order perturbation \(P\Delta P-\rho_\star\Tr(P\Delta P)\).  Once each method has a certified law
\begin{equation}
    \MSE_m(p,R)=b_m(p)^2+\frac{s_m(p)}{R}+\rho_m(p,R),
    \qquad
    \abs{\rho_m(p,R)}\leq \frac{C_m(p)}{R^2},
\end{equation}
method selection becomes a lower-envelope problem with a self-consistency test for every crossing.  This moves the paper from a purely formal Taylor framework toward a mathematically defensible finite-budget theory.

\appendix

\section{Quotient delta-method template}

Let \((U_s,V_s)_{s=1}^B\) be i.i.d.\ with means \((u,v)\), \(v>0\), and covariance matrix
\begin{equation}
    \Sigma=\begin{pmatrix}\sigma_U^2&\sigma_{UV}\\\sigma_{UV}&\sigma_V^2\end{pmatrix}.
\end{equation}
Let \(\widehat u=B^{-1}\sum_sU_s\), \(\widehat v=B^{-1}\sum_sV_s\), and \(\widehat\theta=\widehat u/\widehat v\).  With \(\theta=u/v\),
\begin{align}
    \E\widehat\theta
    &=\theta+\frac{\theta\sigma_V^2-\sigma_{UV}}{Bv^2}+\mathcal R_E(B),\\
    \Var(\widehat\theta)
    &=\frac{\sigma_U^2+\theta^2\sigma_V^2-2\theta\sigma_{UV}}{Bv^2}+\mathcal R_V(B),
\end{align}
where, under the bounded-denominator conditions used in the body, \(\abs{\mathcal R_E(B)}\le A_E/B^2+E_{\rm bad}(B)\) and \(\abs{\mathcal R_V(B)}\le A_V/B^2+V_{\rm bad}(B)\).  This is only the algebraic template behind VD and SV; it is the standard delta-method quotient expansion \citep{VanDerVaart1998}.  The certified versions in \cref{thm:vd_quotient_certified_v4,thm:sv_quotient_law_v2} add denominator concentration and clipping to make the expansion non-asymptotically meaningful.

\section{Notation map}

\begin{center}
\begin{tabular}{ll}
\toprule
Symbol & Meaning\\
\midrule
\(p\) & Physical noise scale\\
\(B\) & Statistical sample count for one method\\
\(R\) & Common physical resource budget after method-specific cost normalization\\
\(\rho_\star\) & Ideal output state\\
\(\rho_p\) & Noisy output state\\
\(\Delta=\mathcal L(\rho_\star)\) & First-order noise perturbation\\
\(O\) & Observable\\
\(\mu_\star\) & Ideal expectation value\\
\(D_M\) & VD denominator \(\Tr(\rho_p^M)\)\\
\(b_m(p)\) & Population residual bias of method \(m\)\\
\(c_m(p)\) & Statistical quotient-bias coefficient\\
\(s_m(p)\) & Resource-normalized sampling coefficient\\
\(a(p)\) & SV acceptance probability\\
\(\delta_{\mathrm{VD}}\) & VD dominant-eigenvector mismatch coefficient\\
\(\delta_{\mathrm{SV}}\) & SV undetectable-sector bias coefficient\\
\bottomrule
\end{tabular}
\end{center}


\begin{thebibliography}{99}

\bibitem[Temme et~al.(2017)]{Temme2017}
K. Temme, S. Bravyi, and J. M. Gambetta.
\newblock Error mitigation for short-depth quantum circuits.
\newblock \emph{Physical Review Letters}, 119:180509, 2017.

\bibitem[Li and Benjamin(2017)]{LiBenjamin2017}
Y. Li and S. C. Benjamin.
\newblock Efficient variational quantum simulator incorporating active error minimization.
\newblock \emph{Physical Review X}, 7:021050, 2017.

\bibitem[Endo et~al.(2018)]{Endo2018}
S. Endo, S. C. Benjamin, and Y. Li.
\newblock Practical quantum error mitigation for near-future applications.
\newblock \emph{Physical Review X}, 8:031027, 2018.

\bibitem[Bonet-Monroig et~al.(2018)]{BonetMonroig2018}
X. Bonet-Monroig, R. Sagastizabal, M. Singh, and T. E. O'Brien.
\newblock Low-cost error mitigation by symmetry verification.
\newblock \emph{Physical Review A}, 98:062339, 2018.

\bibitem[McArdle et~al.(2019)]{McArdle2019}
S. McArdle, X. Yuan, and S. Benjamin.
\newblock Error-mitigated digital quantum simulation.
\newblock \emph{Physical Review Letters}, 122:180501, 2019.

\bibitem[Cai(2021)]{Cai2021}
Z. Cai.
\newblock Quantum error mitigation using symmetry expansion.
\newblock \emph{Quantum}, 5:548, 2021.

\bibitem[Huggins et~al.(2021)]{Huggins2021}
W. J. Huggins, S. McArdle, T. E. O'Brien, J. Lee, N. C. Rubin, S. Boixo, K. B. Whaley, R. Babbush, and J. R. McClean.
\newblock Virtual distillation for quantum error mitigation.
\newblock \emph{Physical Review X}, 11:041036, 2021.

\bibitem[Koczor(2021a)]{Koczor2021PRX}
B. Koczor.
\newblock Exponential error suppression for near-term quantum devices.
\newblock \emph{Physical Review X}, 11:031057, 2021.

\bibitem[Koczor(2021b)]{Koczor2021NJP}
B. Koczor.
\newblock The dominant eigenvector of a noisy quantum state.
\newblock \emph{New Journal of Physics}, 23:123047, 2021.

\bibitem[Czarnik et~al.(2021)]{Czarnik2021}
P. Czarnik, A. Arrasmith, P. J. Coles, and L. Cincio.
\newblock Error mitigation with Clifford quantum-circuit data.
\newblock \emph{Quantum}, 5:592, 2021.

\bibitem[Lowe et~al.(2021)]{Lowe2021}
A. Lowe, M. H. Gordon, P. Czarnik, A. Arrasmith, P. J. Coles, and L. Cincio.
\newblock Unified approach to data-driven quantum error mitigation.
\newblock \emph{Physical Review Research}, 3:033098, 2021.

\bibitem[Kakkar et~al.(2022)]{Kakkar2022}
A. Kakkar, J. Larson, A. Galda, and R. Shaydulin.
\newblock Characterizing error mitigation by symmetry verification in QAOA.
\newblock In \emph{2022 IEEE International Conference on Quantum Computing and Engineering (QCE)}, pages 306--316, 2022.

\bibitem[Bultrini et~al.(2023)]{Bultrini2023}
D. Bultrini, M. H. Gordon, P. Czarnik, A. Arrasmith, M. Cerezo, P. J. Coles, and L. Cincio.
\newblock Unifying and benchmarking state-of-the-art quantum error mitigation techniques.
\newblock \emph{Quantum}, 7:1034, 2023.

\bibitem[Vikstal et~al.(2024)]{Vikstal2024}
P. Vikstal, G. Ferrini, and S. Puri.
\newblock Study of noise in virtual distillation circuits for quantum error mitigation.
\newblock \emph{Quantum}, 8:1441, 2024.

\bibitem[Teo et~al.(2022)]{Teo2022}
Y. S. Teo, S. Shin, H. Kwon, S.-H. Lee, and H. Jeong.
\newblock Virtual distillation with noise dilution.
\newblock arXiv:2210.14753, 2022.

\bibitem[Ijaz et~al.(2024)]{Ijaz2024}
A. Ijaz, C. Huerta Alderete, F. Sauvage, L. Cincio, M. Cerezo, and M. L. Goh.
\newblock More buck-per-shot: Why learning trumps mitigation in noisy quantum sensing.
\newblock arXiv:2410.00197, 2024.

\bibitem[van der Vaart(1998)]{VanDerVaart1998}
A. W. van der Vaart.
\newblock \emph{Asymptotic Statistics}.
\newblock Cambridge University Press, 1998.

\bibitem[Boucheron et~al.(2013)]{Boucheron2013}
S. Boucheron, G. Lugosi, and P. Massart.
\newblock \emph{Concentration Inequalities: A Nonasymptotic Theory of Independence}.
\newblock Oxford University Press, 2013.

\end{thebibliography}


\end{document}
